% Options for packages loaded elsewhere
\PassOptionsToPackage{unicode}{hyperref}
\PassOptionsToPackage{hyphens}{url}
%
\documentclass[
  a4paper,
]{article}
\usepackage{amsmath,amssymb}
\usepackage{iftex}
\ifPDFTeX
  \usepackage[T1]{fontenc}
  \usepackage[utf8]{inputenc}
  \usepackage{textcomp} % provide euro and other symbols
\else % if luatex or xetex
  \usepackage{unicode-math} % this also loads fontspec
  \defaultfontfeatures{Scale=MatchLowercase}
  \defaultfontfeatures[\rmfamily]{Ligatures=TeX,Scale=1}
\fi
\usepackage{lmodern}
% \ifPDFTeX\else
%   % xetex/luatex font selection
%   \setmainfont[]{Libertinus Serif}
%   \setmonofont[Scale=MatchLowercase,Contextuals=Alternate,Ligatures=TeX]{Fira Code}
% \fi
% Use upquote if available, for straight quotes in verbatim environments
\IfFileExists{upquote.sty}{\usepackage{upquote}}{}
\IfFileExists{microtype.sty}{% use microtype if available
  \usepackage[]{microtype}
  \UseMicrotypeSet[protrusion]{basicmath} % disable protrusion for tt fonts
}{}
\makeatletter
\@ifundefined{KOMAClassName}{% if non-KOMA class
  \IfFileExists{parskip.sty}{%
    \usepackage{parskip}
  }{% else
    \setlength{\parindent}{0pt}
    \setlength{\parskip}{6pt plus 2pt minus 1pt}}
}{% if KOMA class
  \KOMAoptions{parskip=half}}
\makeatother
\usepackage{xcolor}
\usepackage[margin=25mm]{geometry}
\usepackage{listings}
\newcommand{\passthrough}[1]{#1}
\lstset{defaultdialect=[5.3]Lua}
\lstset{defaultdialect=[x86masm]Assembler}
\setlength{\emergencystretch}{3em} % prevent overfull lines
\providecommand{\tightlist}{%
  \setlength{\itemsep}{0pt}\setlength{\parskip}{0pt}}
\setcounter{secnumdepth}{5}
\usepackage{xcolor}
\definecolor{codebg}{HTML}{F7F7F7}
\definecolor{codeframe}{HTML}{DDDDDD}

\usepackage{tikz}
\usepackage{tikz-cd}
\usetikzlibrary{trees}
\usetikzlibrary{calc}

\usepackage{bussproofs}
\usepackage{ifthen}
\usepackage{lscape}
\usepackage{graphicx}

\providecommand{\tightlist}{}

\newenvironment{scaledprooftree}[1]
  {\gdef\scalefactor{#1}\begin{center}\proofSkipAmount \leavevmode}%
  {\scalebox{\scalefactor}{\DisplayProof}\proofSkipAmount \end{center}}


\usepackage{listings}
\lstdefinestyle{pandocCode}{
  backgroundcolor=\color{codebg},
  frame=none,                  % il frame lo mette tcolorbox
  framesep=6pt,
  % margini per tenere tutto dentro la box:
  xleftmargin=1.2em,           % ← sposta leggermente a destra TUTTO (incl. numeri)
  xrightmargin=0pt,
  aboveskip=0pt, belowskip=0pt,
  basicstyle=\ttfamily\footnotesize,
  columns=fullflexible,
  keepspaces=true,
  showstringspaces=false,
  breaklines=true,
  breakatwhitespace=true,
  numbers=left,                % ← numeri ON
  numberstyle=\scriptsize\color{gray},
  numbersep=8pt,               % ← spazio tra numeri e codice
  tabsize=2,
}
\lstset{style=pandocCode}

\lstdefinelanguage{lean}{
  morekeywords={theorem,lemma,example,def,inductive,structure,namespace,end,
    intro,intros,assume,suppose,have,show,by,exact,apply,fun,let,match,with},
  sensitive=true,
}

\hbadness=10000

\usepackage[most]{tcolorbox}
\tcbuselibrary{listings, skins, breakable}

% --- TEOREMA (verde stabile) ---
\definecolor{teoremabg}{HTML}{E6F4EC}
\definecolor{teoremaframe}{HTML}{3E8E6D}

% --- PROPOSIZIONE (verde più chiaro) ---
\definecolor{propbg}{HTML}{EDF7F1}
\definecolor{propframe}{HTML}{5FAE8B}

% --- OSSERVAZIONE (verde-grigio neutro) ---
\definecolor{obsbg}{HTML}{EEF3F2}
\definecolor{obsframe}{HTML}{6F8F87}

% --- DEFINIZIONE (azzurro soft ma visibile) ---
\definecolor{defbg}{HTML}{E8F1FB}
\definecolor{defframe}{HTML}{3F7FD6}

% --- ESEMPIO (grigio più deciso) ---
\definecolor{esempbg}{HTML}{F0F0F2}
\definecolor{esempframe}{HTML}{6E6E73}

% --- NOTA (rosso tenue ma visibile) ---
\definecolor{notabg}{HTML}{FBE9E9}
\definecolor{notaframe}{HTML}{C94F4F}

% --- DIMOSTRAZIONE (giallo caldo ma soft) ---
\definecolor{proofbg}{HTML}{FFF6DD}
\definecolor{proofframe}{HTML}{D4A843}


% --- STILE BASE ---
\tcbset{
  mybox/.style={
    enhanced,
    breakable,
    boxrule=0.4pt,
    arc=2mm,
    left=6pt, right=6pt, top=6pt, bottom=6pt,
    width=\textwidth,
    before skip=12pt,
    before=\par\noindent,
    before upper=\setlength{\parskip}{6pt}\setlength{\parindent}{0pt},
    fonttitle=\bfseries,
  }
}

% --- BOX ---

\newtcolorbox{teoremabox}[1][Teorema]{mybox,
  colback=teoremabg,
  colframe=teoremaframe,
  title=#1
}

\newtcolorbox{propositionbox}[1][Proposizione]{mybox,
  colback=propbg,
  colframe=propframe,
  title=#1
}

\newtcolorbox{observationbox}[1][Osservazione]{mybox,
  colback=obsbg,
  colframe=obsframe,
  title=#1
}

\newtcolorbox{definitionbox}[1][Definizione]{mybox,
  colback=defbg,
  colframe=defframe,
  title=#1
}

\newtcolorbox{examplebox}[1][Esempio]{mybox,
  colback=esempbg,
  colframe=esempframe,
  title=#1
}

\newtcolorbox{notebox}[1][Nota]{mybox,
  colback=notabg,
  colframe=notaframe,
  title=#1
}

\newtcolorbox{proofbox}[1][Dimostrazione]{mybox,
  colback=proofbg,
  colframe=proofframe,
  title=#1
}
\ifLuaTeX
  \usepackage{selnolig}  % disable illegal ligatures
\fi
\IfFileExists{bookmark.sty}{\usepackage{bookmark}}{\usepackage{hyperref}}
\IfFileExists{xurl.sty}{\usepackage{xurl}}{} % add URL line breaks if available
\urlstyle{same}
\hypersetup{
  hidelinks,
  pdfcreator={LaTeX via pandoc}}

\author{}
\date{}

\begin{document}

\begin{titlepage}
\thispagestyle{empty}
\begin{center}
{\large\textcolor[HTML]{555555}{Università di Bologna}}\\[0.5em]
{\normalsize\textcolor[HTML]{777777}{Dipartimento di Informatica}}\\[6em]

{\huge\bfseries Algebra e Geometria}\\[0.8em]
{\huge\bfseries Risposte alle domande di teoria}\\[6em]

{\Large Andrea Manzo}\\[1.2em]
{\large A.S.\ 2025 / 2026}\\[6em]
\end{center}
\end{titlepage}


{
\setcounter{tocdepth}{3}
\tableofcontents
}
\newpage

\hypertarget{domande-per-la-prova-di-teoria}{%
\section{Domande per la prova di
teoria}\label{domande-per-la-prova-di-teoria}}

\hypertarget{domanda-1}{%
\subsection{Domanda 1}\label{domanda-1}}

\begin{observationbox}[Domanda]
\begin{enumerate}
\def\labelenumi{\arabic{enumi}.}
\item
  Definizione di vettori linearmente indipendenti. Se possibile, fornire
  esempi dei seguenti vettori, o motivare perché non esistono:

  \begin{itemize}
  \tightlist
  \item
    3 vettori linearmente indipendenti di \(\mathbb{R}^3[x]\);
  \item
    3 vettori linearmente dipendenti di \(\mathbb{R}^3[x]\);
  \item
    5 vettori linearmente dipendenti di \(\mathbb{R}^3[x]\) che non
    generano \(\mathbb{R}^3[x]\);
  \item
    4 vettori linearmente indipendenti di \(\mathbb{R}^3[x]\) che non
    generano \(\mathbb{R}^3[x]\).
  \end{itemize}
\end{enumerate}
\end{observationbox}

\textbf{Definizione:} Sia \(V\) uno spazio vettoriale. I vettori
\(v_1, ..., v_n \in V\) si dicono \textbf{linearmente indipendenti} se
l'unica loro combinazione lineare che dà il vettore nullo è quella con
coefficienti tutti nulli, cioè: \[
a_1v_1 + ... + a_n v_n = \underline{0} \quad\Longrightarrow\quad a_1 = a_2 = ... = a_n = 0
\]

\textbf{Esempi:}

\begin{itemize}
\tightlist
\item
  3 vettori linearmente indipendenti di \(\mathbb{R}^3[x]\): \[
    \begin{aligned}
    &p_1 = 1 + x \\
    &p_2 = 1 + x + x^2 \\
    &p_3 = 1 + x + x^2 + x^3
    \end{aligned}
    \]
\item
  3 vettori linearmente dipendenti di \(\mathbb{R}^3[x]\): \[
    \begin{aligned}
    &p_1 = 1 + x \\
    &p_2 = 2 + 2x \\
    &p_3 = 3 + 3x 
    \end{aligned}
    \]
\item
  5 vettori linearmente dipendenti di \(\mathbb{R}^3[x]\) che non
  generano \(\mathbb{R}^3[x]\): \[
    \begin{aligned}
    &p_1 = 1 + x \\
    &p_2 = 2 + 2x \\
    &p_3 = 3 + 3x \\
    &p_4 = 1 + x^2 \\
    &p_5 = 2 + 2x^2
    \end{aligned}
    \]
\item
  4 vettori linearmente indipendenti di \(\mathbb{R}^3[x]\) che non
  generano \(\mathbb{R}^3[x]\): Impossibile in quanto
  \(\mathbb{R}^3[x]\) ha dimensione 4 e perciò, per GEL, 4 vettori
  linearmente indipendenti costituiscono una base di \(\mathbb{R}^3[x]\)
  e quindi generano.
\end{itemize}

\hypertarget{domanda-2}{%
\subsection{Domanda 2}\label{domanda-2}}

\begin{observationbox}[Domanda]
\begin{enumerate}
\def\labelenumi{\arabic{enumi}.}
\setcounter{enumi}{1}
\tightlist
\item
  Definizione di vettori linearmente indipendenti e di combinazione
  lineare. Enunciare e dimostrare il teorema: i vettori
  \(v_1,\ldots,v_n\) di uno spazio vettoriale \(V\) sono linearmente
  dipendenti se e solo se\ldots{}
\end{enumerate}
\end{observationbox}

\textbf{Definizione:} Sia \(V\) uno spazio vettoriale, siano
\(v_1, v_2, \dots, v_n \in V\). Una \textbf{combinazione lineare} di
\(v_1, \dots, v_n\) è un vettore \(v\) del tipo: \[
v = a_1v_1 + a_2v_2 + \dots + a_nv_n
\] con \(a_1, \dots ,a_n \in \mathbb{R}\)

\textbf{Teorema:} I vettori \(v_1, ..., v_n \in V\) sono dipendenti
\(\iff\) uno di essi è \textbf{combinazione lineare degli altri}.

\textbf{Dimostrazione:}

Dividiamo la dimostrazione:

\(\Rightarrow\) : Per ipotesi \(v_1, ..., v_n\) sono dipendenti, quindi
esiste una combinazione lineare \[
a_1v_1 + ... + a_n v_n = \underline{0} 
\] tale che gli \(a_i\) non sono tutti nulli. Sia \(a_k \neq 0\), allora
\[
a_k v_k = - a_1 v_1 -a_2 v_2 - ... - a_{k-1}v_{k-1}-a_{k+1}v_{k+1} -...- a_n v_n
\] Dividiamo entrambi i membri per \(a_k\): \[
v_k = - \frac{a_1}{a_k} v_1 - ... - \frac{a_{k-1}}{a_k}v_{k-1} -... - \frac{a_n}{a_k}v_n
\] Vediamo quindi che \(v_k\) è combinazione lineare di
\(v_1, ..., v_{k-1}, v_{k+1}, ..., v_n\)

\(\Leftarrow\) : Per ipotesi un vettore \(v_k\) è combinazione lineare
degli altri, quindi \[
v_k = a_1 v_1 + a_2 v_2 + ... + a_{k-1}v_{k-1}+a_{k+1}v_{k+1} +...+ a_n v_n
\] Se aggiungiamo \(-v_k\) a entrambi i membri: \[
a_1 v_1 + a_2 v_2 + ... + a_{k-1}v_{k-1}-v_k +a_{k+1}v_{k+1} +...+ a_n v_n = \underline{0}
\] con \(-1 \neq 0\) (il coefficiente di \(v_k\)). I vettori sono
dipendenti per definizione.

\hypertarget{domanda-3}{%
\subsection{Domanda 3}\label{domanda-3}}

\begin{observationbox}[Domanda]
\begin{enumerate}
\def\labelenumi{\arabic{enumi}.}
\setcounter{enumi}{2}
\tightlist
\item
  Definizione di \(\langle v_1,\ldots,v_n\rangle\). Dimostrare che
  \(\langle v_1,\ldots,v_n\rangle\) è un sottospazio.
\end{enumerate}
\end{observationbox}

\textbf{Definizione:} Dati i vettori \(v_1,\ldots,v_n\) indichiamo con
\(\langle v_1,\ldots,v_n\rangle\) l'insieme di tutte le combinazioni
lineari dei vettori \(v_1,\ldots,v_n\) , ovvero: \[
\langle v_1, \dots, v_n \rangle = \{a_1v_1 + \dots + a_n v_n \mid a_1, \dots, a_n \in \mathbb{R}\}
\]

\(\langle v_1,\ldots,v_n\rangle\) è detto anche \textbf{sottospazio
generato} (o span). Infatti dimostriamo ora che è un sottospazio.

\textbf{Dimostrazione:}

Mostriamo che \(\underline{0} \in \langle v_1,\ldots,v_n\rangle\):
poiché \[
\underline{0} = 0 v_1 + 0 v_2 + \dots + 0 v_n
\] è combinazione lineare di \(v_1, \dots, v_n\) allora
\(\underline{0} \in \langle v_1,\ldots,v_n\rangle\)

Mostriamo che è chiuso rispetto alla somma: sia
\(x = a_1 v_1 + \dots + a_n v_n\) e sia
\(y = b_1 v_1 + \dots + b_n v_n\) . Allora: \[
x + y = a_1 v_1 + \dots + a_n v_ +  b_1 v_1 + \dots + b_n v_n = (a_1 + b_1)v_1 + (a_2 + b_2) v_2 + \dots + (a_n + b_n) v_n
\] che è ancora una combinazione lineare di \(v_1, \dots, v_n\) perciò
\(x + y \in \langle v_1,\ldots,v_n\rangle\).

Mostriamo che è chiuso rispetto al prodotto per scalare: sia
\(x = a_1 v_1 + \dots + a_n v_n\), ,sia \(\lambda \in \mathbb{R}\).
Allora: \[
\lambda x = \lambda (a_1 v_1 + \dots + a_n v_n) = \lambda a_1 v_1 + \dots + \lambda a_n v_n
\] che è ancora una combinazione lineare di \(v_1, \dots, v_n\) perciò
\(\lambda x \in \langle v_1,\ldots,v_n\rangle\).

\hypertarget{domanda-4}{%
\subsection{Domanda 4}\label{domanda-4}}

\begin{observationbox}[Domanda]
\begin{enumerate}
\def\labelenumi{\arabic{enumi}.}
\setcounter{enumi}{3}
\tightlist
\item
  Definizione di base di un sottospazio. Esempio di una base di
  \(\mathbb{R}^3[x]\) diversa dalla base canonica. Dimostrare che uno
  spazio vettoriale finitamente generato ha sempre una base.
\end{enumerate}
\end{observationbox}

\textbf{Definizione:} Sia \(V\) uno spazio vettoriale. L'insieme
\(\{v_1, \dots v_n\}\) è una \textbf{base} di \(V\) se:

\begin{enumerate}
\def\labelenumi{\arabic{enumi}.}
\tightlist
\item
  \(V = \langle v_1, \dots , v_n\rangle\), ovvero se \(v_1, \dots, v_n\)
  generano \(V\)
\item
  \(v_1, \dots, v_n\) sono linearmente indipendenti
\end{enumerate}

\textbf{Esempio:} una base di \(\mathbb{R}^3[x]\) diversa dalla base
canonica è \(\beta = \{ 3x^3, x^2 + 1, 2x, 5\}\)

\bigskip

Mostriamo che uno spazio vettoriale finitamente generato ha sempre una
base.

\textbf{Dimostrazione:}

Sia \(V = \langle v_1, ..., v_n \rangle\). Se \(V = \{\underline 0\}\)
allora una base di \(V\) è \(\emptyset\), altrimenti consideriamo
\(v_1, ..., v_n\):

\begin{itemize}
\tightlist
\item
  se sono indipendenti abbiamo finito e \(\beta = \{v_1, ..., v_n\}\) è
  la base cercata
\item
  se sono dipendenti invece sappiamo che uno di essi e combinazione
  lineare degli altri, sia esso \(v_i\). Ma allora: \[
  V = \langle v_1, ..., v_n \rangle = V = \langle v_1, ..., v_{i-1}, v_{i+1}, ...,  v_n \rangle
  \] abbiamo cioè rimosso \(v_i\). Ripartiamo con il procedimento.
\end{itemize}

Dopo un numero finito di passi avremo finito e i vettori che rimangono
sono base di \(V\).

\hypertarget{domanda-5}{%
\subsection{Domanda 5}\label{domanda-5}}

\begin{observationbox}[Domanda]
\begin{enumerate}
\def\labelenumi{\arabic{enumi}.}
\setcounter{enumi}{4}
\tightlist
\item
  Definizione di dimensione di uno spazio vettoriale.
\end{enumerate}
\end{observationbox}

\textbf{Definizione:} La \textbf{dimensione} di uno spazio vettoriale è
il numero di elementi di una sua base. Tale valore è unico in quanto
ogni base ha lo stesso numero di elementi.

\hypertarget{domanda-6}{%
\subsection{Domanda 6}\label{domanda-6}}

\begin{observationbox}[Domanda]
\begin{enumerate}
\def\labelenumi{\arabic{enumi}.}
\setcounter{enumi}{5}
\tightlist
\item
  Coordinate di un vettore rispetto ad una base ordinata: definizione,
  esistenza e unicità.
\end{enumerate}
\end{observationbox}

\textbf{Definizione (e proposizione di esistenza e unicità):} Sia
\(\beta=\{v_1,\ldots,v_n\}\) una base ordinata per lo spazio vettoriale
\(V\) e sia \(v\in V\). Allora esiste ed è unica la \(n\)-upla di
scalari \((\lambda_1,\ldots,\lambda_n)\) tale che\\
\[  
v=\lambda_1 v_1+\cdots+\lambda_n v_n
\] Gli scalari \((\lambda_1,\ldots,\lambda_n)\) si dicono le
\textbf{coordinate} di \(v\) rispetto alla base ordinata \(\beta\).
Scriviamo: \[
(v)_{\beta} = (\lambda_1,\ldots,\lambda_n)
\]

\textbf{Dimostrazione:} Mostriamo esistenza e unicità delle coordinate:

\begin{itemize}
\item
  Esistenza: \(\beta=\{v_1,\ldots,v_n\}\) è base di \(V\), quindi
  \(V = \langle v_1 ,\dots ,v_n\rangle\), quindi \(v\) è combinazione
  lineare di \(v_1 ,\dots ,v_n\) cioè esistono
  \(\lambda_1,\ldots,\lambda_n\) tali che
  \(v=\lambda_1 v_1+\cdots+\lambda_n v_n\).
\item
  Unicità: Sia \[
  v=\lambda_1 v_1+\cdots+\lambda_n v_n = \gamma_1 v_1+\cdots+\gamma_n v_n
  \] Quindi: \[
  \lambda_1 v_1+\cdots+\lambda_n v_n - \gamma_1 v_1-\cdots-\gamma_n v_n = \underline{0}
  \] perciò, raccogliendo: \[
    (\lambda_1 - \gamma_1)v_1 + ... + (\lambda_n - \gamma_n)v_n = \underline{0}
   \] Quindi, siccome \(v_1,\ldots,v_n\) sono indipendenti, per
  definizione: \[
    (\lambda_1 - \gamma_1) = (\lambda_2 - \gamma_2) = ... = (\lambda_n - \gamma_n) = \underline{0}
   \] Perciò: \[
   \lambda_1 = \gamma_1, \quad \lambda_2 = \gamma_2, \quad ... , \quad \lambda_n = \gamma_n
   \]
\end{itemize}

\hypertarget{domanda-7}{%
\subsection{Domanda 7}\label{domanda-7}}

\begin{observationbox}[Domanda]
\begin{enumerate}
\def\labelenumi{\arabic{enumi}.}
\setcounter{enumi}{6}
\tightlist
\item
  Sia \(U=\langle v_1,\ldots,v_k\rangle\subseteq \mathbb{R}^n\). Perché
  l'uso dell'algoritmo di Gauss in modo diretto permette di trovare una
  base di \(U\)?
\end{enumerate}
\end{observationbox}

L'uso dell'algoritmo di Gauss in modo diretto permette di trovare una
base di \(U\) grazie a due proposizioni:

\begin{enumerate}
\def\labelenumi{\arabic{enumi}.}
\tightlist
\item
  Le righe non nulle di una matrice a scala sono linearmente
  indipendenti.
\item
  Sia \(A \in M_{m \times n}(\mathbb{R})\) , allora le operazioni
  elementari sulle righe di \(A\) non cambiano il sottospazio generato
  dalle righe di \(A\) .
\end{enumerate}

Quindi possiamo costruire una matrice \(A\) che abbia come righe i
vettori \(v_1,\ldots,v_k\) ed eseguire poi l'algoritmo di Gauss su di
essa. Alla fine le righe non nulle della matrice a scala \(A'\) saranno
indipendenti ed essendo che non cambia il sottospazio generato, generano
sempre \(U\). Essendo indipendenti e generando, per definizione sono una
base di \(U\).

\hypertarget{domanda-8}{%
\subsection{Domanda 8}\label{domanda-8}}

\begin{observationbox}[Domanda]
\begin{enumerate}
\def\labelenumi{\arabic{enumi}.}
\setcounter{enumi}{7}
\tightlist
\item
  Siano \(v_1,\ldots,v_k\in \mathbb{R}^n\). Perché l'uso dell'algoritmo
  di Gauss in modo diretto permette stabilire se \(v_1,\ldots,v_k\) sono
  linearmente indipendenti?
\end{enumerate}
\end{observationbox}

L'uso dell'algoritmo di Gauss in modo diretto permette di stabilire se
dei vettori sono linearmente indipendenti grazie a due proposizioni:

\begin{enumerate}
\def\labelenumi{\arabic{enumi}.}
\tightlist
\item
  Le righe non nulle di una matrice a scala sono linearmente
  indipendenti.
\item
  Sia \(A \in M_{m \times n}(\mathbb{R})\) , allora le operazioni
  elementari sulle righe di \(A\) non cambiano il sottospazio generato
  dalle righe di \(A\) .
\end{enumerate}

Sia \(V\) il sottospazio generato da \(v_1,\ldots,v_k\). Possiamo
costruire una matrice \(A\) che abbia come righe i vettori
\(v_1,\ldots,v_k\) ed eseguire poi l'algoritmo di Gauss su di essa. Se
alla fine le righe non nulle della matrice a scala \(A'\) sono sempre
\(k\) , vuol dire che i vettori \(v_1,\ldots,v_k\) sono indipendenti, in
quanto le \(k\) righe sono indipendenti e generano sempre \(V\) e perciò
sono una base di \(V\) e perciò \(V\) ha dimensione \(k\) e perciò
essendo \(v_1,\ldots,v_k\) generatori, sono una base e quindi
indipendenti.

\hypertarget{domanda-9}{%
\subsection{Domanda 9}\label{domanda-9}}

\begin{observationbox}[Domanda]
\begin{enumerate}
\def\labelenumi{\arabic{enumi}.}
\setcounter{enumi}{8}
\tightlist
\item
  Equazioni cartesiane di un sottospazio \(W\) di \(\mathbb{R}^n\) e
  motivazione teorica del loro calcolo a partire da un insieme di
  generatori di \(W\).
\end{enumerate}
\end{observationbox}

\textbf{Definizione:} Sia \(W\) il sottospazio tale che: \[
W=\{x\in\mathbb{R}^n \mid A\underline{x}=\underline{0}\}
\] per una qualche matrice \(A\in M_{k\times n}(\mathbb{R})\). Allora,
le equazioni del sistema lineare omogeneo
\(A\underline{x}=\underline{0}\) prendono il nome di \textbf{equazioni
cartesiane} di \(W\).

\textbf{Calcolo delle equazioni cartesiane:}

Siano \(v_1,\dots,v_m\in\mathbb{R}^n\) tali che
\(W=\langle v_1,\dots,v_m\rangle\) . Allora \[
W=
\{
\lambda_1 v_1 +\cdots+\lambda_m v_m
\mid
\lambda_1,\dots,\lambda_m \in\mathbb{R}
\}
\]

Sia \(B\in M_{n\times m}(\mathbb{R})\) la matrice che ha come colonne i
vettori \(v_1,\dots,v_m\). Un vettore \[
x=(x_1,\dots,x_n)\in\mathbb{R}^n
\] appartiene a \(V\) se e solo se esistono
\(\lambda_1,\dots,\lambda_m\in\mathbb{R}\) tali che \[
x=\lambda_1v_1+\cdots+\lambda_mv_m
\] Equivalentemente, esiste \[
\lambda=(\lambda_1,\dots,\lambda_m)\in\mathbb{R}^m
\] tale che \[
B\lambda=x
\] Per il teorema di Rouché-Capelli, ciò equivale a richiedere che: \[
\operatorname{rk}(B)=\operatorname{rk}(B|x)
\]

Applicando l'algoritmo di Gauss alla matrice aumentata \((B|x)\), si
ottiene una matrice della forma

\[
\begin{pmatrix}
\bullet & \cdots & \cdots & \;\;|\;\cdots \\
0 & \bullet & \cdots & \;\;|\;\cdots \\
0 & 0 & \bullet & \;\;|\;\cdots \\
0 & 0 & 0 & |\;\diamond \\
0 & 0 & 0 & |\;\diamond
\end{pmatrix}
\]

dove:

\begin{itemize}
\tightlist
\item
  i simboli \(\bullet\) indicano i pivot di riga;
\item
  i simboli \(\diamond\) indicano espressioni lineari nelle variabili
  \(x_1,\dots,x_n\).
\end{itemize}

Per avere \(\operatorname{rk}(B)=\operatorname{rk}(B|x)\) è sufficiente
imporre che tutte le espressioni indicate con \(\diamond\) siano nulle.

Le equazioni così ottenute sono le equazioni cartesiane del sottospazio
\(W\).

\hypertarget{domanda-10}{%
\subsection{Domanda 10}\label{domanda-10}}

\begin{observationbox}[Domanda]
\begin{enumerate}
\def\labelenumi{\arabic{enumi}.}
\setcounter{enumi}{9}
\tightlist
\item
  Definizione di applicazione lineare e di nucleo. Dimostrare che una
  applicazione lineare è iniettiva se e solo se il suo nucleo contiene
  solo il vettore nullo.
\end{enumerate}
\end{observationbox}

\textbf{Definizione:} Siano \(V, W\) due spazi vettoriali, sia
\(F: V \to W\) una funzione. \(F\) si dice \textbf{applicazione lineare}
se:

\begin{enumerate}
\def\labelenumi{\arabic{enumi}.}
\tightlist
\item
  \(F(v_1 + v_2) = F(v_1) + F(v_2)\)
\item
  \(F(\lambda v) = \lambda F(v)\)
\end{enumerate}

\textbf{Definizione:} Sia \(F: V \to W\) una applicazione lineare,
chiamiamo \textbf{nucleo} o \textbf{kernel} di \(F\) l'insieme: \[
\text{Ker}(F) = \{v \in V \mid F(v) = \underline{0}\}
\]

\textbf{Dimostrazione:}

Mostriamo che \(F\) è iniettiva
\(\iff \text{Ker}(F) = \{\underline{0}\}\):

\(\Rightarrow\) : Per ipotesi \(F\) è iniettiva, quindi
\(F(u) = F(v) \Longrightarrow u = v, \; \forall u, v \in V\) .

Sia \(v \in \text{Ker}(F)\), perciò per definizione
\(F(v) = \underline{0}\) . Inoltre, essendo \(F\) lineare, sappiamo
anche che \(F(\underline{0}) = \underline{0}\), ma allora, per
definizione di iniettività, \(v = \underline{0}\), perciò
\(\text{Ker}(F) = \{\underline{0}\}\)

\(\Leftarrow\): Per ipotesi \(\text{Ker}(F) = \{\underline{0}\}\) .
Dobbiamo mostrare che
\(F(u) = F(v) \Longrightarrow u = v, \; \forall u, v \in V\). Notiamo
che: \[
F(u) = F(v) \Longrightarrow F(u) - F(v) = \underline{0}  \Longrightarrow F(u - v) = \underline{0}
\] Ma allora per definizione di kernel e per le ipotesi: \[
u -v \in \text{Ker}(F) \Longrightarrow u - v = \underline{0} \Longrightarrow u = v
\]

\hypertarget{domanda-11}{%
\subsection{Domanda 11}\label{domanda-11}}

\begin{observationbox}[Domanda]
\begin{enumerate}
\def\labelenumi{\arabic{enumi}.}
\setcounter{enumi}{10}
\tightlist
\item
  Definizione di applicazione lineare e immagine. Completare e
  dimostrare l'enunciato: se \(F:V\to W\) è un'applicazione lineare e
  \(B=\{\ldots\}\) è una base di \ldots, allora
  \(\text{Im}(F)=\langle\ldots\rangle\).
\end{enumerate}
\end{observationbox}

\textbf{Definizione:} Siano \(V, W\) due spazi vettoriali, sia
\(F: V \to W\) una funzione. \(F\) si dice \textbf{applicazione lineare}
se:

\begin{enumerate}
\def\labelenumi{\arabic{enumi}.}
\tightlist
\item
  \(F(v_1 + v_2) = F(v_1) + F(v_2)\)
\item
  \(F(\lambda v) = \lambda F(v)\)
\end{enumerate}

\textbf{Definizione:} Sia \(F: V \to W\) una applicazione lineare, si
chiama \textbf{immagine} di \(F\) l'insieme: \[
\text{Im}(F) = \{F(v) \mid v \in V\} = \{w \in W \mid w = F(v), v \in V\}
\]

\textbf{Teorema:} se \(F:V\to W\) è un'applicazione lineare e
\(\beta=\{ v_1, \dots, v_n\}\) è una base di \(V\) allora
\(\text{Im}(F)=\langle F(v_1), \dots ,F(v_n)\rangle\)

\textbf{Dimostrazione:}

Sappiamo che: \[
\text{Im}(L) = \{L(v) \mid v \in V\} = S_1
\] e che: \[
\langle L(v_1), \dots , L(v_n) \rangle = \{\lambda_1 L(v_1) + \dots + \lambda_n L(v_n) \mid (\lambda_1, \dots \lambda_n \in \mathbb{R}\} = S_2
\] Vogliamo mostrare che \(S_1 = S_2\).

\begin{itemize}
\item
  Mostriamo che \(S_1 \subseteq S_2\) :

  Sia \(w \in S_1\) , perciò \(w = L(v)\) con
  \(v \in V = \langle v_1, \dots, v_n \rangle\) ,
  \(v = \lambda_1 v_1 + \dots + \lambda_n v_n\) . Allora: \[
   w = L(v) = L(\lambda_1 v_1 + \dots + \lambda_n v_n) = \lambda_1 L(v_1) + \dots + \lambda_n L(v_n)
   \] Perciò \(w \in \langle L(v_1), \dots , L(v_n) \rangle = S_2\) .
\item
  Mostriamo che \(S_2 \subseteq S_1\):

  Sia \(w \in S_2\), perciò
  \(w = \lambda_1 L(v_1) + \dots + \lambda_n L(v_n)\) . Allora: \[
   w = \lambda_1 L(v_1) + \dots + \lambda_n L(v_n) = L(\lambda_1 v_1 + \dots + \lambda_n v_n)
   \] Perciò \(w \in \text{Im}(L) = S_1\).
\end{itemize}

Quindi abbiamo dimostrato \(S_1 = S_2\) .

\hypertarget{domanda-12}{%
\subsection{Domanda 12}\label{domanda-12}}

\begin{observationbox}[Domanda]
\begin{enumerate}
\def\labelenumi{\arabic{enumi}.}
\setcounter{enumi}{11}
\tightlist
\item
  Data una matrice \(A\in M_{m,n}(\mathbb{R}^n)\), definire
  l'applicazione lineare \(L_A\) ad essa associata e dare la motivazione
  teorica del calcolo del suo nucleo.
\end{enumerate}
\end{observationbox}

\textbf{Definizione:} sia \(A \in M_{m\times n}(\mathbb{R})\) ,
definiamo l'applicazione lineare \(L_A\): \[
L_A : \mathbb{R}^n \to \mathbb{R}^m
\] \[
\quad\underline{x} \to A\underline{x}
\] dove \(\underline{x}\) è un \textbf{vettore colonna}. Quindi
l'immagine di un vettore \(v\) tramite \(L_A\) è il prodotto righe per
colonne di \(A\) con \(v\).

\(L_A\) può essere definita dando le immagini di una base del dominio
(base canonica, o facendo poi un cambio di base), e in particolare le
colonne di \(A\) saranno \(C_1 = L(e_1), \dots, C_n = L(e_n)\) ; oppure
può essere definita dando la legge
\(L(x_1, x_2, \dots, x_n) = (y_1, y_2, \dots, y_m)\) .

\textbf{Calcolo del nucleo}:

Il nucleo di \(L_A\) è
\(\text{Ker}(L_A) = \{v \in \mathbb{R}^n \mid L_A(v) = \underline{0}\} = \{v \in \mathbb{R}^n \mid A\cdot v = \underline{0}\}\)
. Perciò per calcolarlo ci basta risolvere il sistema associato
\(Av = \underline{0}\) . Se il sistema ammette una sola soluzione essa è
quella nulla e perciò \(\text{Ker}(L_A) = \{\underline{0}\}\) ,
altrimenti avremo infinite soluzioni (e quindi infiniti vettori nel
nucleo) dipendenti da \(n - r\) parametri, dove \(r\) è il rango righe
della matrice \(A\). Il kernel avrà quindi dimensione \(n - r\) .

\hypertarget{domanda-13}{%
\subsection{Domanda 13}\label{domanda-13}}

\begin{observationbox}[Domanda]
\begin{enumerate}
\def\labelenumi{\arabic{enumi}.}
\setcounter{enumi}{12}
\tightlist
\item
  Data una matrice \(A\in M_{m,n}(\mathbb{R}^n)\), definire
  l'applicazione lineare \(L_A\) ad essa associata e dare la motivazione
  teorica del calcolo della sua immagine.
\end{enumerate}
\end{observationbox}

\textbf{Definizione:} sia \(A \in M_{m\times n}(\mathbb{R})\) ,
definiamo l'applicazione lineare \(L_A\): \[
L_A : \mathbb{R}^n \to \mathbb{R}^m
\] \[
\quad\underline{x} \to A\underline{x}
\] dove \(\underline{x}\) è un \textbf{vettore colonna}. Quindi
l'immagine di un vettore \(v\) tramite \(L_A\) è il prodotto righe per
colonne di \(A\) con \(v\).

\(L_A\) può essere definita dando le immagini di una base del dominio
(base canonica, o facendo poi un cambio di base), e in particolare le
colonne di \(A\) saranno \(C_1 = L(e_1), \dots, C_n = L(e_n)\) ; oppure
può essere definita dando la legge
\(L(x_1, x_2, \dots x_n) = (y_1, y_2, \dots y_m)\) .

\textbf{Calcolo dell'immagine}:

Sappiamo che \(\text{Im}(L) = \langle L(e_1), \dots ,L(e_n) \rangle\)
perciò le colonne di \(A\) generano l'immagine. Per trovarne una base
possiamo mettere le colonne di \(A\) in riga e applicare Gauss in modo
diretto per trovare quali di esse sono linearmente indipendenti. Questo
equivale ad applicare Gauss su \(A^T\). L'immagine di \(A\) ha quindi
dimensione uguale al rango colonne di \(A\) .

\hypertarget{domanda-14}{%
\subsection{Domanda 14}\label{domanda-14}}

\begin{observationbox}[Domanda]
\begin{enumerate}
\def\labelenumi{\arabic{enumi}.}
\setcounter{enumi}{13}
\tightlist
\item
  Dimostrare il teorema: siano \(V,W\) spazi vettoriali, sia
  \(B=\{v_1,\ldots,v_n\}\) una base di \(V\) e siano
  \(w_1,\ldots,w_n\in W\). Allora esiste una ed una sola applicazione
  lineare \(F:V\to W\) tale che \(F(v_1)=w_1,\ldots,F(v_n)=w_n\).
\end{enumerate}
\end{observationbox}

\textbf{Dimostrazione:}

Definiamo innanzitutto \(F\), cioè assegniamo una legge che ad ogni
\(v \in V\) associa un elemento di \(W\).

Sia \(v \in V\) scriviamo \(v = \lambda_1 v_1 + \dots + \lambda_n v_n\)
. Poniamo \[
F(v) = \lambda_1 w_1 + \dots + \lambda_n w_n
\]

Ora mostriamo che \(F\) è lineare.

Siano \(u, v \in V\) mostriamo che \(F(u + v) = F(u) + F(v)\) : \[
u = a_1 v_1 + \dots + a_n v_n \qquad v = b_1 v_1 + \dots + b_n v_n
\] \[
F(u) = a_1 w_1 + \dots + a_n w_n \qquad F(v) = b_1 w_1 + \dots + b_n w_n
\] Allora calcoliamo \(F(u + v)\): \[
\begin{aligned}
F(u + v) &= F(a_1 v_1 + \dots + a_n v_n + b_1 v_1 + \dots + b_n v_n) \\
&= F((a_1 + b_1) v_1 + \dots + (a_n + b_n) v_n) \\
&= (a_1 + b_1) w_1 + \dots + (a_n + b_n) w_n \\
&= a_1 w_1 + \dots + a_n w_n + b_1 w_1 + \dots + b_n w_n \\
&= F(u) + F(v)
\end{aligned}
\]

Sia \(u\in V\) mostriamo che \(F(\lambda u) = \lambda F(u)\) \[
\begin{aligned}
F(\lambda u) &= F(\lambda a_1 v_1 + \lambda a_2 v_2 + \dots + \lambda a_n v_n ) \\
&= \lambda a_1 w_1 + \lambda a_2 w_2 + \dots + \lambda a_n w_n \\
&= \lambda (a_1 w_1 + \dots + a_n w_n) \\
&= \lambda F(u)
\end{aligned}
\] Abbiamo quindi dimostrato che \(F\) è lineare.

Dobbiamo ora mostrare che \(F\) è unica:

Sia \(G: V \to W\) lineare, tale che \(G(v_i) = w_i\) . Mostriamo che
\(F = G\):

Sia \(v \in V\) scriviamo \(v = \lambda_1 v_1 + \dots + \lambda_n v_n\)
. Calcoliamo \(G(v)\): \[
\begin{aligned}
G(v) &= G(\lambda_1 v_1 + \dots + \lambda_n v_n) \\
&= \lambda_1 G(v_1) + \dots + \lambda_n G(v_n)\\
&= \lambda_1 w_1 + \dots + \lambda_n w_n \\
&= F(v)
\end{aligned}
\] Quindi abbiamo dimostrato che \(F\) è unica.

\hypertarget{domanda-15}{%
\subsection{Domanda 15}\label{domanda-15}}

\begin{observationbox}[Domanda]
\begin{enumerate}
\def\labelenumi{\arabic{enumi}.}
\setcounter{enumi}{14}
\tightlist
\item
  Enunciare e dimostrare il Teorema della Dimensione.
\end{enumerate}
\end{observationbox}

\textbf{Teorema:} Sia \(F: V \to W\) lineare. Allora:\\
\[
\dim(V) = \dim(\text{Ker}(F)) + \dim(\text{Im}(F))
\]

\textbf{Dimostrazione:}

Sia \(\dim(V) = n, \; \dim(\text{Ker}(F)) = r\). Mostriamo che
\(\dim(\text{Im}(F)) = n - r\) .

Sia \(\{v_1, \dots, v_r \}\) base di \(\text{Ker}(F)\).
\(v_1, \dots, v_r\) sono linearmente indipendenti quindi posso
completare \(\{v_1, \dots, v_r \}\) a una base di \(V\) (aggiungendo
\(n-r\) vettori indipendenti): \[
\beta = \{v_1, .., v_r, z_{r+1},..., z_{n}\}
\] Mostriamo che \(\{F(z_{r+1}), ..., F(z_n)\}\) è una base di
\(\text{Im}(F)\) , mostrando così che che \(\dim(\text{Im}(F)) = n - r\)
.

Sappiamo che: \[
\text{Im}(F) = \langle F(v_1), \dots , F(v_r), F(z_{r+1}), ..., F(z_n) \rangle = \langle F(z_{r+1}), ..., F(z_n)\rangle
\]poiché \(F(v_1) = F(v_2) = \dots = F(v_r) = \underline{0}\) in quanto
fanno pare del kernel.

Mostriamo che \(F(z_{r+1}), ..., F(z_n)\) sono linearmente indipendenti.

Sia \[
\lambda_{r+1} F(z_{r+1})+ ...+ \lambda_{n}F(z_n) = \underline{0}
\] dobbiamo mostrare che
\(\lambda_{z+1} = ... = \lambda_n = \underline{0}\) (definizione di
lineare indipendenza).

Essendo \(F\) lineare,
\(F(\lambda_{r+1}z_{r+1} + ...+ \lambda_{n} z_n) = \underline{0}\) .
Perciò dato \[
z = \lambda_{r+1}z_{r+1} + ...+ \lambda_{n} z_n
\] vale che: \[
F(z) = \underline{0} \quad\Longrightarrow \quad z \in \text{Ker}(F) = \langle v_1, \dots,  v_r \rangle
\] Perciò \[
z = \lambda_1 v_1 + ... + \lambda_r v_r
\] Perciò essendo che ovviamente \(z - z = \underline{0}\): \[
\lambda_1 v_1 + ... + \lambda_r v_r - \lambda_{r+1}z_{r+1} - ...- \lambda_{n} z_n = \underline{0}
\]

Essendo \(\beta\) linearmente indipendente deve valere che \[
\lambda_1 = \lambda_2 = ... = -\lambda_{z+1} = ... = - \lambda_n = \underline{0}
\]

In particolare quindi \[
\lambda_{z+1} = ... = \lambda_n = \underline{0}
\] Quindi \(F(z_{r+1}), ..., F(z_n)\) sono linearmente indipendenti.
Quindi sono base di \(\text{Im}(F)\). Quindi vale che
\(\dim(\text{Im}(F)) = n - r\) .

\hypertarget{domanda-16}{%
\subsection{Domanda 16}\label{domanda-16}}

\begin{observationbox}[Domanda]
\begin{enumerate}
\def\labelenumi{\arabic{enumi}.}
\setcounter{enumi}{15}
\tightlist
\item
  Definire il rango righe e il rango colonne di una matrice e dimostrare
  che sono uguali.
\end{enumerate}
\end{observationbox}

\textbf{Definizione:} Sia \(A \in M_{m \times n}(\mathbb{R})\).

Il \textbf{rango righe} di \(A\) è \(rr(A) =\dim\) del sottospazio
generato dalle righe di \(A\).

Il \textbf{rango colonne} di \(A\) è \(rc(A) = \dim\) del sottospazio
generato dalle colonne di \(A\).

\textbf{Dimostrazione:}

Dimostriamo che \(rr(A) = rc(A)\) . Consideriamo l'applicazione lineare
\(L_A : \mathbb{R}^n \to \mathbb{R}^m\) tale che
\(L_A (\underline{x}) = A\underline{x}\) .

Per il teorema della dimensione: \[
n = \dim(\text{Ker} (L_A)) + \dim(\text{Im} (L_A))
\] \(\text{Ker}(L_A)\) è l'insieme delle soluzioni del sistema
\(A\underline{x} = \underline{0}\) e ha dimensione \(n - rr(A)\) perché
dopo aver ridotto \(A\) a scala, abbiamo \(rr(A)\) pivot e le soluzioni
dipendono \(n - rr(A)\) parametri.

\(\text{Im}(L_A) = \langle L_A(v_1), \dots , L_A(v_n) \rangle\) =
sottospazio generato dalle colonne di \(A\), e ha \(\dim = rc(A)\)\\
Otteniamo: \[
n = (n -rr(A)) + rc(A) \quad \Rightarrow \quad rr(A) = rc(A)
\]

\hypertarget{domanda-17}{%
\subsection{Domanda 17}\label{domanda-17}}

\begin{observationbox}[Domanda]
\begin{enumerate}
\def\labelenumi{\arabic{enumi}.}
\setcounter{enumi}{16}
\tightlist
\item
  Definizione di isomorfismo tra spazi vettoriali e di spazi vettoriali
  isomorfi. Enunciare e possibilmente dimostrare il teorema che
  caratterizza gli spazi vettoriali isomorfi.
\end{enumerate}
\end{observationbox}

\textbf{Definizione:} Sia \(F : V \to W\) una applicazione lineare.
\(F\) si i dice \textbf{isomorfismo} se è iniettiva e suriettiva (ovvero
biettiva).

\textbf{Definizione:} Due spazi vettoriali \(V, W\) si dicono
\textbf{isomorfi} se esiste un isomorfismo \(F: V \to W\)

\textbf{Teorema:} Due spazi vettoriali \(V, W\) sono isomorfi \(\iff\)
hanno la stessa dimensione.

\textbf{Dimostrazione:}

\(\Rightarrow\): Per ipotesi \(V\) è isomorfo a \(W\) e sia
\(F: V \to W\) isomorfismo. Per il teorema della dimensione:

\begin{itemize}
\tightlist
\item
  \(F\) iniettiva \(\Rightarrow \text{Ker}(F)\) ha dimensione 0
\item
  \(F\) suriettiva \(\Rightarrow \text{Im}(F) = W\)
\end{itemize}

Quindi \(\dim(V) = 0 + \dim(W)\). Quindi \(V\) e \(W\) hanno la stessa
dimensione.

\(\Leftarrow\): Per ipotesi \(\dim(V) = \dim(W) = n\) . Sia
\(\beta = \{v_1 ... v_n\}\) base di \(V\) e sia
\(\tilde{\beta} = \{w_1 ... w_n\}\) base di \(W\).

Per il teorema di esistenza delle applicazioni lineari esiste un'
(unica) applicazione lineare \(F: V \to W\) tale che
\(F(v_1) = w_1, ..., F(v_n) = w_n\).

Mostriamo che \(F\) è un isomorfismo:

Sappiamo che
\(\text{Im}(F) = \langle F(v_1) .... F(v_n) \rangle = \langle w_1, ..., w_n \rangle = W\),
quindi \(F\) è suriettiva.\\
Per il teorema della dimensione
\(\dim(V) = \dim(\text{Ker}(F)) + \dim(\text{Im}(F))\), ma allora,
poiché \(\dim(V) = \dim(\text{Im}(F)) = n\), vale che
\(\dim(\text{Ker}(F)) = 0\), perciò \(F\) è iniettiva.

\hypertarget{domanda-18}{%
\subsection{Domanda 18}\label{domanda-18}}

\begin{observationbox}[Domanda]
\begin{enumerate}
\def\labelenumi{\arabic{enumi}.}
\setcounter{enumi}{17}
\tightlist
\item
  Controimmagine (o preimmagine) di un vettore tramite una applicazione
  lineare: definizione e caratterizzazione.
\end{enumerate}
\end{observationbox}

\textbf{Definizione:} Sia \(F: A \to B\) una funzione. Sia \(b \in B\).
Allora: \[
F^{-1}(b) = \{a \in A \mid F(a) = b\} \subseteq A
\] è la \textbf{controimmagine} o \textbf{preimmagine} o
\textbf{antiimmagine} di un \(b\) tramite la funzione \(F\).

\textbf{Teorema di struttura della controimmagine:} Sia \(F : V \to W\)
una applicazione lineare e sia \(w \in W\). Allora innanzitutto: \[
w \in \text{Im}(F) \iff F^{-1}(w) \not = \emptyset
\] Se \(w \in \text{Im}(F)\), fissiamo un qualsiasi \(v \in F^{-1}(w)\).
Allora: \[
F^{-1}(w) = \{v + z \mid z \in \text{Ker}(F) \}
\] Le controimmagini sono quindi dei traslati di \(\text{Ker}(F)\).

\textbf{Dimostrazione:}

Per definizione: \[
F^{-1}(w) = \{u \in V \mid F(u) = w\}
\] Sia \(v \in F^{-1}(w)\). Ora, per mostrare che \[
F^{-1}(w) = \{v + z \mid z \in \text{Ker}(F) \}
\] mostriamo la doppia inclusione tra insiemi. Iniziamo mostrando che \[
\{v + z \mid z \in \text{Ker}(F) \} \subseteq F^{-1}(w) 
\] Sia \(u = v + z\) con \(z \in \text{Ker}(F)\). Allora: \[
F(u) = F(v + z) = F(v) + F(z) = w + 0 = w
\] Quindi \(u \in F^{-1}(w)\) .

Dobbiamo ora mostrare che \[
F^{-1}(w) \subseteq \{v + z \mid z \in \text{Ker}(F) \}
\] Perciò sia \(u \in F^{-1}(w)\), dobbiamo mostrare che possiamo
scrivere \(u\) come \(u = v + z\), con \(z \in \text{Ker}(F)\).
Mostriamo che vale \(z = u - v\). Notiamo che \[
F(z) = F(u - v) = F(u) - F(v) = w - w = \underline{0}
\] che è corretto in quanto vogliamo che \(z \in \text{Ker}(F)\). Perciò
\(z = u - v\). Perciò possiamo ridurci a dimostrare \(u = v + (u - v)\),
che ovviamente si semplifica in \(u = u\). Ovvio.

\hypertarget{domanda-19}{%
\subsection{Domanda 19}\label{domanda-19}}

\begin{observationbox}[Domanda]
\begin{enumerate}
\def\labelenumi{\arabic{enumi}.}
\setcounter{enumi}{18}
\tightlist
\item
  Enunciare il teorema di Rouché-Capelli e dimostrarne la prima parte.
\end{enumerate}
\end{observationbox}

\textbf{Teorema:} Sia \(A\underline{x} = \underline{b}\) un sistema
lineare di \(m\) equazioni in \(n\) incognite. Esso ha soluzione
\(\iff r(A) = r(A \mid b)\).

Inoltre se \(r(A) = r(A \mid b) = r\), il sistema ha:

\begin{itemize}
\tightlist
\item
  1 sola soluzione se \(r=n\)
\item
  Infinite soluzioni che dipendono da \(n-r\) parametri se \(r < n\)
\end{itemize}

\textbf{Dimostrazione prima parte:}

Sia \(L_A : \mathbb{R}^n \to \mathbb{R}^m\) l'applicazione lineare
associata ad \(A\), cioè \(L_A(\underline{x}) = A \underline{x}\) .

Consideriamo questa catena di coimplicazioni (\(\iff\)): \[
\begin{aligned}
\text{Il sistema } A\underline{x} = \underline{b} \text{ ha soluzione }\iff& \\
L_A^{-1}(\underline{b}) \not = \emptyset \iff&\\
\underline{b} \in \text{Im}(L_A) \iff&\\
\underline{b} \in \langle L_A (e_1), \dots, L_A(e_n) \rangle \iff& \\
\underline{b} \in \langle \text{colonne di } A \rangle
\end{aligned}
\] ma sappiamo che se
\(\underline{b} \in \langle \text{colonne di } A \rangle\) allora: \[
\langle \text{colonne di } A \rangle = \langle \underline{b}, \text{colonne di } A \rangle
\] perciò: \[
\begin{aligned}
\iff& \dim(\langle \underline{b},  \text{colonne di } A \rangle)  =\dim(\langle \text{colonne di } A \rangle) \\
\iff& r(A \mid \underline{b}) = r(A)
\end{aligned}
\]

\hypertarget{domanda-20}{%
\subsection{Domanda 20}\label{domanda-20}}

\begin{observationbox}[Domanda]
\begin{enumerate}
\def\labelenumi{\arabic{enumi}.}
\setcounter{enumi}{19}
\tightlist
\item
  Determinante di una matrice: definizione e proprietà.
\end{enumerate}
\end{observationbox}

\textbf{Definizione:} Il \textbf{determinante} è una funzione che
associa ad ogni matrice quadrata \(A\) di ordine \(n\) un numero reale,
indicato con \[
\det(A)
\] tale che valgano le seguenti proprietà:

\begin{enumerate}
\def\labelenumi{\arabic{enumi}.}
\item
  Se la riga \(j\)-esima di \(A\) è somma di due vettori \(u\) e \(v\),
  allora il determinante di \(A\) è la somma dei determinanti delle
  matrici ottenute sostituendo alla riga \(j\)-esima rispettivamente
  \(u\) e \(v\).
\item
  Se la riga \(j\)-esima di \(A\) è del tipo \(\lambda u\), con
  \(\lambda \in \mathbb{R}\), allora \[
  \det(A) = \lambda \det(B)
  \]

  dove \(B\) è la matrice ottenuta sostituendo la \(j\)-esima riga con
  \(u\).
\item
  Se due righe di \(A\) sono uguali, allora \[
  \det(A)=0
  \]
\item
  Vale che \[
  \det(I)=1
  \]
\end{enumerate}

\textbf{Proprietà del determinante:}

Siano \(A\) e \(B\) matrici quadrate di ordine \(n\).

\begin{enumerate}
\def\labelenumi{\arabic{enumi}.}
\tightlist
\item
  Se \(B\) si ottiene da \(A\) scambiando due righe, allora \[
  \det(B) = -\det(A)
  \]
\item
  Se \(B\) si ottiene sommando ad una riga di \(A\) una combinazione
  lineare delle altre righe, allora \[
  \det(B)=\det(A)
  \]
\item
  Se \(A\) è una matrice triangolare superiore oppure triangolare
  inferiore, cioè i coefficienti al di sotto (rispettivamente al di
  sopra) della diagonale principale sono tutti uguali a zero, allora il
  suo determinante è il prodotto degli elementi della diagonale
  principale: \[
  \det(A)=a_{11}a_{22}\cdots a_{nn}
  \]
\end{enumerate}

\hypertarget{domanda-21}{%
\subsection{Domanda 21}\label{domanda-21}}

\begin{observationbox}[Domanda]
\begin{enumerate}
\def\labelenumi{\arabic{enumi}.}
\setcounter{enumi}{20}
\tightlist
\item
  Definizione di matrice quadrata invertibile. Se una matrice quadrata è
  invertibile allora il suo determinante è diverso da zero.
\end{enumerate}
\end{observationbox}

\textbf{Definizione:} Una matrice \(A \in M_{n\times n}(\mathbb{R})\) si
dice \textbf{invertibile} se esiste una matrice
\(B \in M_{n\times n}(\mathbb{R})\), tale che: \[
AB = BA = I
\]La matrice \(B\) si dice l'\textbf{inversa} di \(A\) e si indica con
\(A^{-1}\).

\textbf{Dimostrazione:}

Mostriamo che una matrice \(A \in M_{n}(\mathbb{R})\) è invertibile
\(\Longrightarrow \det(A) \not = 0\) :

Sia \(A\) invertibile e sia \(B\) l'inversa di \(A\). Quindi \(AB = I\)
. Per Binet vale che: \[
\det(AB) = \det(A)\det(B) = \det(I) = 1
\] quindi \[
\det(A)\det(B)=1 \Rightarrow \det(A) \not = 0
\]

\hypertarget{domanda-22}{%
\subsection{Domanda 22}\label{domanda-22}}

\begin{observationbox}[Domanda]
\begin{enumerate}
\def\labelenumi{\arabic{enumi}.}
\setcounter{enumi}{21}
\tightlist
\item
  Definizione di matrice quadrata invertibile. Se \(A,B\) sono matrici
  quadrate tali che \(AB=I\), allora \(B\) è l'inversa di \(A\).
\end{enumerate}
\end{observationbox}

\textbf{Definizione:} Una matrice \(A \in M_{n\times n}(\mathbb{R})\) si
dice \textbf{invertibile} se esiste una matrice
\(B \in M_{n\times n}(\mathbb{R})\), tale che: \[
AB = BA = I
\]La matrice \(B\) si dice l'\textbf{inversa} di \(A\) e si indica con
\(A^{-1}\).

\textbf{Dimostrazione:}

Mostriamo che se \(A,B\) sono matrici quadrate tali che \(AB=I\), allora
\(B\) è l'inversa di \(A\).

Innanzitutto sappiamo che: \[
AB = I \;\Longrightarrow\; \det(AB) = \det(I) = 1
\] Inoltre per Binet: \[
\;\Longrightarrow\; \det(A)\cdot\det(B) = 1 \;\Longrightarrow\; \det(A) \not = 0
\] quindi \(A\) è invertibile, quindi esiste \(A^{-1}\).

Moltiplichiamo a sinistra entrambi i membri di \(AB = I\) per
\(A^{-1}\): \[
A^{-1}AB = A^{-1}I \;\Longrightarrow\; (A^{-1}A)B = A^{-1} \;\Longrightarrow\; IB = A^{-1} \;\Longrightarrow\; B = A^{-1}
\]

\hypertarget{domanda-23}{%
\subsection{Domanda 23}\label{domanda-23}}

\begin{observationbox}[Domanda]
\begin{enumerate}
\def\labelenumi{\arabic{enumi}.}
\setcounter{enumi}{22}
\tightlist
\item
  Enunciare il seguente teorema e dimostrare qualche implicazione
  (''teoremone''). Sia \(F=L_A\) un endomorfismo di \(\mathbb{R}^n\):
  sono equivalenti: \((a)\) \(F\) è iniettiva, \((b)\) \ldots{}
\end{enumerate}
\end{observationbox}

\textbf{Teorema:} Sia \(F : \mathbb{R}^n \to \mathbb{R}^n\) una
applicazione lineare e sia \(A\) la matrice associata a \(F\) rispetto
alla base canonica. Le seguenti affermazioni sono equivalenti:

\begin{enumerate}
\def\labelenumi{\arabic{enumi}.}
\tightlist
\item
  \(F\) è un isomorfismo
\item
  \(F\) è iniettiva
\item
  \(F\) è suriettiva
\item
  \(\dim(\text{Im}(F))=n\)
\item
  \(r(A) = n\)
\item
  Le colonne di \(A\) sono linearmente indipendenti
\item
  Le righe di \(A\) sono linearmente indipendenti
\item
  Il sistema \(A\underline{x} = \underline{0}\) ha un'unica soluzione
\item
  \(\forall\, b \in \mathbb{R}^n\) il sistema
  \(A\underline{x} = \underline{b}\) ha un'unica soluzione
\item
  \(A\) è invertibile
\item
  \(\det(A) \not = 0\)
\end{enumerate}

\textbf{Dimostrazione:}

\begin{itemize}
\tightlist
\item
  1 \(\Rightarrow\) 2 per definizione di isomorfismo
\item
  2 \(\Rightarrow\) 3 per il teorema della dimensione, poiché se \(F\) è
  iniettiva, \(\dim(\text{Ker}(F)) = 0\) quindi
  \(\dim(\text{Im}(F)) = n\), quindi \(F\) è suriettiva
\item
  3 \(\iff\) 4, infatti \(F\) è suriettiva \(\iff \dim(\text{Im}(F))=n\)
\item
  4 \(\Rightarrow\) 1 poiché \(4 \Rightarrow 3\) e inoltre per il
  teorema della dimensione
  \(\dim(\text{Ker}(F)) = n - \dim(\text{Im}(F)) = n - n =0\), quindi
  \(F\) è iniettiva, quindi \(F\) è isomorfismo
\item
  4 \(\Rightarrow\) 5,
  \(\dim(\text{Im}(F))=n \Rightarrow \dim(\langle F(e_1), ..., F(e_n) \rangle) = n\)
  ed è uguale alla dimensione del sottospazio generato dalle colonne,
  che è proprio \(r(A)\)
\item
  5 \(\Rightarrow\) 6:
  \(r(A) = n \Rightarrow \dim(\langle \text{colonne di A} \rangle) = n \Rightarrow \langle \text{colonne di A} \rangle = \mathbb{R}^n \Rightarrow\)
  per GEL sono linearmente indipendenti
\item
  6 \(\Rightarrow\) 4: le colonne ci \(A\) sono linearmente indipendenti
  , quindi \(\dim(\langle \text{colonne di A} \rangle) = n\). Ma
  \(\text{Im}(F)\) è proprio il sottospazio generato dalle colonne,
  quindi \(\dim(\text{Im}(F))=n\)
\item
  5 \(\iff\) 7 ovvio per GEL
\item
  8 \(\iff\) 2, perché \(F\) è iniettiva
  \(\iff \text{Ker}(F) = \{\underline{0}\} \iff \{\text{soluzione di } A\underline{x} = \underline{0}\} = \{\underline{0}\}\)
\item
  9 \(\iff\) 5, per il teorema di Rouché-Capelli
\item
  10 \(\iff\) 11 per la proposizione sul calcolo dell'inversa,
  dimostrata precedentemente.
\item
  1 \(\iff\) 10, \(F\) è biunivoca \(\iff\) \(F\) ha un'inversa \(\iff\)
  \(A\) ammette un'inversa
\end{itemize}

\hypertarget{domanda-24}{%
\subsection{Domanda 24}\label{domanda-24}}

\begin{observationbox}[Domanda]
\begin{enumerate}
\def\labelenumi{\arabic{enumi}.}
\setcounter{enumi}{23}
\tightlist
\item
  Matrice associata ad una applicazione lineare rispetto a due basi
  fissate \(\beta\) del dominio e \(\widetilde{\beta}\) del codominio:
  come si costruisce e perché permette di determinare l'immagine di un
  qualsiasi vettore. La matrice
  \(I_{\beta \mathcal{C}} = M_{\beta}^{\mathcal{C}}(\text{id})\)
  associata all'identità di \(\mathbb{R}^n\) .
\end{enumerate}
\end{observationbox}

Siano

\begin{itemize}
\tightlist
\item
  \(F : V \to W\) applicazione lineare,
\item
  \(\beta = \{v_1, \dots, v_n\}\) base ordinata di \(V\)
\item
  \(\widetilde{\beta} = \{w_1, \dots, w_m\}\) base ordinata di \(W\)
\end{itemize}

Allora esiste una matrice \[
A = A_{\beta \widetilde{\beta}} = M^{\widetilde{\beta}}_{\beta}(F) \quad \in M_{m\times n}(\mathbb{R})
\] tale che se conosco \(A\) conosco \(F\), ovvero: \[
(F(v))_{\widetilde{\beta}} = A\cdot(v)_{\beta}
\] Inoltre nell' \(i\)-esima colonna di \(A\) ci sono le coordinate di
\(F(v_i)\) rispetto alla base \(\widetilde{\beta}\) : \[
A=M^{\widetilde{\beta}}_{\beta}(F)  
=  
\left(  
\begin{array}{cccc}  
\big| & \big| & & \big| \\  
(F(v_1))_{\widetilde{\beta}}  
&  
(F(v_2))_{\widetilde{\beta}}  
&  
\cdots  
&  
(F(v_n))_{\widetilde{\beta}}  
\\  
\big| & \big| & & \big|  
\end{array}  
\right)
\] Perciò, dato un vettore \(v\), se conosco \(A\) posso calcolare
\(F(v)\) : \[
v \to \text{coordinate } (v)_{\beta} \to A \cdot (v)_\beta = (F(v))_{\widetilde{\beta}} \to \text{dalle coordinate rispetto a } \widetilde{\beta} \text{ ottengo } F(v)
\] In pratica le coordinate di \(F(v)\) rispetto a \(\widetilde{\beta}\)
sono il prodotto righe per colonne di \(A\) per le coordinate di \(v\)
rispetto alla base di \(V\).

\bigskip

La matrice
\(I_{\beta \mathcal{C}} = M_{\beta}^{\mathcal{C}}(\text{id})\) associata
all'identità di \(\mathbb{R}^n\) è la matrice che ha per colonne i
vettori \(v_1, \dots , v_n\) della base ordinata \(\beta\) .

\hypertarget{domanda-25}{%
\subsection{Domanda 25}\label{domanda-25}}

\begin{observationbox}[Domanda]
\begin{enumerate}
\def\labelenumi{\arabic{enumi}.}
\setcounter{enumi}{24}
\tightlist
\item
  Matrice associata ad una applicazione lineare rispetto a due basi
  fissate \(\beta\) del dominio e \(\widetilde{\beta}\) del codominio:
  come si costruisce e perché permette di determinare l'immagine di un
  qualsiasi vettore. La matrice
  \(I_{\beta \beta} = M_{\beta}^{\beta}(\text{id})\) associata
  all'identità di \(\mathbb{R}^n\) .
\end{enumerate}
\end{observationbox}

Prima parte come domanda 24.

La matrice \(I_{\beta \beta} = M_{\beta}^{\beta}(\text{id})\) associata
all'identità di \(\mathbb{R}^n\) è la matrice identità (\(I\)), in
quanto per un generico vettore della base \(v_i\), nella colonna i-esima
di \(I_{\beta \beta}\) ci sono le coordinate di \(v_i\) rispetto a
\(\beta\), ma essendo \(v_1 \in \beta\) le coordinate saranno
\((0, 0, \dots , 1, \dots, 0)\), ovvero tutti zero, tranne 1 nella
posizione \(i\)-esima.

\hypertarget{domanda-26}{%
\subsection{Domanda 26}\label{domanda-26}}

\begin{observationbox}[Domanda]
\begin{enumerate}
\def\labelenumi{\arabic{enumi}.}
\setcounter{enumi}{25}
\tightlist
\item
  Dimostrare che \(I_{\mathcal{C} \beta} = I_{\beta \mathcal{C}}^{-1}\)
  , cioè che \(M^B_C(\text{id})=(M^C_B(\text{id}))^{-1}\).
\end{enumerate}
\end{observationbox}

Consideriamo la seguenti applicazioni lineari di identità:

\begin{itemize}
\tightlist
\item
  \(\text{id}' : \mathbb{R}^n (\text{base } \mathcal{C}) \to \mathbb{R}^n (\text{base } \beta)\)
  . La matrice associata è \(I_{\mathcal{C} \beta}\) .
\item
  \(\text{id}'' : \mathbb{R}^n (\text{base } \beta) \to \mathbb{R}^n (\text{base } \mathcal{C})\)
  . La matrice associata è \(I_{\beta \mathcal{C}}\) .
\end{itemize}

E consideriamo la funzione composta \(\text{id}' \circ \text{id}''\) .
Essendo \(\text{id}' ,\text{id}''\) lineari, la matrice associata alla
composta \(\text{id}' \circ \text{id}''\) è
\(I_{\beta \mathcal{C}} \cdot I_{\mathcal{C} \beta}\) (la matrice
dell'applicazione lineare che agisce prima si trova più a destra). Lo
schema relativo è il seguente:

\begin{center}
\begin{tikzcd}[column sep=huge,row sep=tiny]
\mathbb{R}^n
\arrow[r, "\mathrm{id}'", "I_{\mathcal{C} \beta }"']
\arrow[rr, bend right=35, dashed]
&
\mathbb{R}^n
\arrow[r, "\mathrm{id}''", "I_{\beta \mathcal{C} }"']
&
\mathbb{R}^n
\\
{\scriptstyle \mathcal{C}}
&
{\scriptstyle \beta}
&
{\scriptstyle \mathcal{C}}
\end{tikzcd}
\end{center}

Essendo \(\text{id}' \circ \text{id}'' = \text{id}\) , la matrice
associata alla composta è
\(I_{\beta \mathcal{C}} \cdot I_{\mathcal{C} \beta} = I_{\mathcal{C} \mathcal{C}} = I\)
. Quindi: \[
I_{\beta \mathcal{C}} \cdot I_{\mathcal{C} \beta} = I \;\Longrightarrow\; I_{\mathcal{C} \beta} = I_{\beta \mathcal{C}}^{-1}
\]

\hypertarget{domanda-27}{%
\subsection{Domanda 27}\label{domanda-27}}

\begin{observationbox}[Domanda]
\begin{enumerate}
\def\labelenumi{\arabic{enumi}.}
\setcounter{enumi}{26}
\tightlist
\item
  Definizione di endomorfismo diagonalizzabile e di matrice
  diagonalizzabile. Dimostrare che un endomorfismo è diagonalizzabile se
  e solo se la matrice associata è diagonalizzabile.
\end{enumerate}
\end{observationbox}

\textbf{Definizione:} Un endomorfismo \(F: V \to V\) si dice
\textbf{diagonalizzabile} se esiste una base \(\beta\) di \(V\) tale che
\(M^{{\beta}}_{\beta}(F)\) sia diagonale.

\textbf{Definizione:} Una matrice quadrata \(A\) si dice
\textbf{diagonalizzabile} se è simile ad una matrice diagonale, cioè se
esiste una matrice invertibile \(P\) tale che \(P^{-1}AP = D\) con \(D\)
diagonale.

\textbf{Dimostrazione:}

\(\Rightarrow\): Sia \(F\) diagonalizzabile, allora esiste \(\beta\)
base di \(\mathbb{R}^n\) tale che \(M^{{\beta}}_{\beta}(F) = D\) è
diagonale. La formula del cambio di base dice che
\(D = I^{{-1}}_{\beta C}AI_{\beta C}\) quindi \(A\) è simile a una
matrice diagonale, quindi \(A\) è diagonalizzabile.

\(\Leftarrow\): Supponiamo \(A\) diagonalizzabile, cioè esiste \(P\)
invertibile tale che \(P^{-1}AP = D\) con \(D\) diagonale. Vogliamo
trovare una base \(\beta\) tale che \(M^{{\beta}}_{\beta}(F)\) sia
diagonale.

Osserviamo che siccome \(P\) è invertibile, le colonne di \(P\) sono
linearmente indipendenti, quindi formano una base di \(\mathbb{R}^n\).
Sia \(\beta = \{v_1, ..., v_n\)\}, dove \(v_1, ..., v_n\) sono le
colonne di \(P\).

Abbiamo che \(I_{\beta \mathcal{C}} = P\) . Sappiamo per la formula del
cambio di base che \[
M^{{\beta}}_{\beta}(F) = I^{{-1}}_{\beta C}AI_{\beta C} = P^{-1}AP = D
\]

\hypertarget{domanda-28}{%
\subsection{Domanda 28}\label{domanda-28}}

\begin{observationbox}[Domanda]
\begin{enumerate}
\def\labelenumi{\arabic{enumi}.}
\setcounter{enumi}{27}
\tightlist
\item
  Definizione di endomorfismo diagonalizzabile e di autovettore di un
  endomorfismo. Dimostrare che se \(F\) è un endomorfismo di
  \(\mathbb{R}^n\) ed esiste una base di \(\mathbb{R}^n\) costituita da
  autovettori di \(F\) allora \(F\) è diagonalizzabile.
\end{enumerate}
\end{observationbox}

\textbf{Definizione:} Un endomorfismo \(F: V \to V\) si dice
\textbf{diagonalizzabile} se esiste una base \(\beta\) di \(V\) tale che
\(M^{{\beta}}_{\beta}(F)\) sia diagonale.

\textbf{Definizione:} Sia \(F: V \to V\) un endomorfismo. Un vettore
\(v \in V\) si dice \textbf{autovettore} di \(F\) se: \[
v \not = \underline{0} \;\;\text{ e }\;\; F(v) = \lambda v \,,\; \lambda \in \mathbb{R}
\] \(\lambda\) si dice \textbf{autovalore} dell'autovettore \(v\).

\textbf{Dimostrazione:}

Sia \(\beta = \{v_1, ... v_n\}\) una base di \(V\) costituita da
autovettori. Mostriamo che \(M^{{\beta}}_{\beta}(F)\) è diagonale.
Essendo \(v_1, ..., v_n\) autovettori, \(F(v_i) = \lambda_i v_i\) .
Cerchiamo ora \((F(v_i))_{\beta}\), osserviamo che: \[
F(v_i) = \lambda_i v_i = 0 v_1 + ... + \lambda_i v_i + ... + 0 v_n
\] e quindi: \[
(F(v_i))_{\beta} = (0 , ... , \lambda_i , ... , 0)
\] Perciò: \[
M^{{\beta}}_{\beta}(F) =
\begin{pmatrix}
\lambda_1 & 0 & \cdots & 0\\
0 & \lambda_2 & \cdots & 0\\
\vdots & \vdots & \ddots & \vdots\\
0 & 0 & \cdots & \lambda_n
\end{pmatrix}
\]

Notiamo che tale matrici è diagonale e che inoltre che sulla diagonale
ci sono gli autovalori ordinati.

\hypertarget{domanda-29}{%
\subsection{Domanda 29}\label{domanda-29}}

\begin{observationbox}[Domanda]
\begin{enumerate}
\def\labelenumi{\arabic{enumi}.}
\setcounter{enumi}{28}
\tightlist
\item
  Definizione di endomorfismo diagonalizzabile e di autovettore di un
  endomorfismo. Dimostrare che se un endomorfismo \(F\) di
  \(\mathbb{R}^n\) è diagonalizzabile allora esiste una base di
  \(\mathbb{R}^n\) costituita da autovettori di \(F\).
\end{enumerate}
\end{observationbox}

Prima parte come domanda 28.

\textbf{Dimostrazione:}

Supponiamo che \(F\) sia diagonalizzabile, allora esiste una base
\(\beta\) tale che \(M^{{\beta}}_{\beta}(F)\) è diagonale, ovvero
abbiamo sulla diagonale principale i valori
\(\lambda_1, ..., \lambda_n\). Mostriamo che \(\beta = \{v_1, ...v_n\}\)
è costituita da autovettori. Sappiamo che sulla \(i\)-esima colonna di
\(M^{{\beta}}_{\beta}(F)\) ci sono le coordinate di \(F(v_i)\) rispetto
a \(\beta\), quindi: \[
F(v_i) = 0 v_1 + ... + \lambda_i v_i + ... + 0 v_n = \lambda_i v_i
\] quindi \(v_i\) è autovettore di autovalore \(\lambda_i\) .

\hypertarget{domanda-30}{%
\subsection{Domanda 30}\label{domanda-30}}

\begin{observationbox}[Domanda]
\begin{enumerate}
\def\labelenumi{\arabic{enumi}.}
\setcounter{enumi}{29}
\tightlist
\item
  Il polinomio caratteristico di una matrice: definizione e sue
  proprietà.
\end{enumerate}
\end{observationbox}

\textbf{Definizione:} Se \(A \in M_n(\mathbb{R})\), sia \(A - xI\) la
matrice: \[
A -
\begin{pmatrix}
x & 0 & \cdots & 0\\
0 & x & \cdots & 0\\
\vdots & \vdots & \ddots & \vdots\\
0 & 0 & \cdots & x
\end{pmatrix}
= 
\begin{pmatrix}
a_{11}-x & a_{12} & \cdots & a_{1n}\\
a_{21} & a_{22}-x & \cdots & a_{2n}\\
\vdots & \vdots & \ddots & \vdots\\
a_{n1} & a_{n2} & \cdots & a_{nn}-x
\end{pmatrix}
\] Il \textbf{polinomio caratteristico} di \(A\) è: \[
P_A(x) = \det(A - xI)
\]

\bigskip

\textbf{Proprietà del polinomio caratteristico:}

\begin{enumerate}
\def\labelenumi{\arabic{enumi}.}
\tightlist
\item
  Ha grado \(n\) e termine direttivo \((-1)^n x^n\)
\item
  Il termine noto è \(\det(A)\)
\item
  Il termine di grado 1 è \(-(a_{11} + a_{22} + ... + a_{nn})\)
\end{enumerate}

\hypertarget{domanda-31}{%
\subsection{Domanda 31}\label{domanda-31}}

\begin{observationbox}[Domanda]
\begin{enumerate}
\def\labelenumi{\arabic{enumi}.}
\setcounter{enumi}{30}
\tightlist
\item
  Relazione di similitudine tra matrici. Matrici simili hanno lo stesso
  determinante? e lo stesso polinomio caratteristico? Motivare con la
  dimostrazione (se fatta a lezione).
\end{enumerate}
\end{observationbox}

\textbf{Definizione:} Siano \(A, B \in M_n(\mathbb{R})\) , \(A\) è
\textbf{simile} a \(B\) se esiste una matrice invertibile
\(P \in M_n(\mathbb{R})\), tale che \[
A = P^{-1}BP
\]

La relazione di similitudine tra matrici è una \textbf{relazione di
equivalenza}:

\begin{enumerate}
\def\labelenumi{\arabic{enumi}.}
\tightlist
\item
  Riflessiva: \(A\) è simile ad \(A\), poiché \(A = I^{-1}AI\)
\item
  Simmetrica: se \(A\) è simile a \(B\), allora \(A = P^{-1}BP\), allora
  \(PAP^{-1} = P(P^{-1}BP)P^{-1} = B\), che posso scrivere come
  \(B = (P^{-1})^{-1}AP^{-1}\)
\item
  Transitiva: \(A\) simile a \(B\), B simile a \(C\) \(\Rightarrow\)
  \(A\) simile a \(C\). (Dimostrazione omessa)
\end{enumerate}

\textbf{Proposizione:} Matrici simili hanno lo stesso determinante.

\textbf{Dimostrazione:}

Per ipotesi \(A = P^{-1}BP\). Calcoliamo:

\[
\det(A) = \det(P^{-1}BP)
\] per Binet: \[
= \det(P^{-1})\cdot \det(B) \cdot \det(P) = \det(B) \cdot \det(P^{-1}P) = \det(B) \cdot \det(I) = \det(B)
\]

\bigskip

\textbf{Proposizione:} Matrici simili hanno lo stesso polinomio
caratteristico

\textbf{Dimostrazione:}

Per ipotesi \(A = P^{-1}BP\). Calcoliamo:

\[ 
\begin{aligned}
P_A(x) &=\det(A - xI) = \det(P^{-1}BP - xI) = \\
&=\det(P^{-1}BP - P^{-1}PxI) = \det(P^{-1}BP - P^{-1}(xI)P) = \\
&=\det(P^{-1}(B-xI)P) =
\end{aligned}
\] ora, per binet: \[
\begin{aligned}
&=\det(P^{-1})\det(B-xI) \det(P) = P_B(x)\det(P^{-1})\det(P)\\
&=P_B(x)\det(P^{-1}P) =P_B(x)\det(I) = \\
&=P_B(x)
\end{aligned}
\]

\hypertarget{domanda-32}{%
\subsection{Domanda 32}\label{domanda-32}}

\begin{observationbox}[Domanda]
\begin{enumerate}
\def\labelenumi{\arabic{enumi}.}
\setcounter{enumi}{31}
\tightlist
\item
  Esempio di due matrici che hanno lo stesso polinomio caratteristico ma
  non sono simili (spiegando il motivo).
\end{enumerate}
\end{observationbox}

Consideriamo le matrici: \[
A = 
\begin{pmatrix}
1 & 0 \\
0 & 1 \\
\end{pmatrix}  
\qquad
B = 
\begin{pmatrix}
1 & 1 \\
0 & 1 \\
\end{pmatrix}  
\] Calcoliamone i polinomi caratteristici: \[
P_A(x) = \det (A - xI) = \det
\begin{pmatrix}
1 -x & 0 \\
0 & 1 -x\\
\end{pmatrix}  
= (1-x)^2
\] \[
P_B(x) = \det (B - xI) = \det
\begin{pmatrix}
1 -x & 1 \\
0 & 1 -x\\
\end{pmatrix}  
= (1-x)^2
\] Hanno lo stesso polinomi o caratteristico. Mostriamo ora che non sono
simili.

Notiamo che \(A\) è diagonalizzabile (in quanto è diagonale). Mostriamo
che invece \(B\) non lo è. Gli autovalori di \(B\) sono: \[
\lambda = 1 \qquad m_a(1) = 2
\] Calcoliamo \(m_g(1)\). Calcoliamo quindi la dimensione
dell'autospazio \(V_1\): \[
V_1 = \text{Ker}(B - I) = \text{Ker} \begin{pmatrix}
0 & 1 \\
0 & 0 \\
\end{pmatrix}
= \{(t, 0) \mid t \in \mathbb{R}\} = \langle(1, 0) \rangle
\] Quindi \(\text{dim}(V_1) = 1\), quindi \(m_g(1) = 1 \neq m_a(1)\) ,
quindi \(B\) non è diagonalizzabile.

Ed essendo \(A\) diagonalizzabile e \(B\) no, \(A\) e \(B\) non possono
essere simili in quando la similitudine è una relazione transitiva:
\(A\) è simile a una matrice \(D\) diagonale, ma se \(B\) è simile ad
\(A\), anche \(B\) dovrebbe essere simile a \(D\) (impossibile perché
\(B\) non è diagonalizzabile).

\hypertarget{domanda-33}{%
\subsection{Domanda 33}\label{domanda-33}}

\begin{observationbox}[Domanda]
\begin{enumerate}
\def\labelenumi{\arabic{enumi}.}
\setcounter{enumi}{32}
\tightlist
\item
  Definizione di autovalore di un endomorfismo, molteplicità algebrica e
  geometrica. Proprietà equivalenti al fatto che un endomorfismo o una
  matrice siano diagonalizzabili (solo enunciato).
\end{enumerate}
\end{observationbox}

\textbf{Definizione:} Sia \(F: V \to V\) un endomorfismo. Un vettore
\(v \in V\) si dice \textbf{autovettore} di \(F\) se: \[
v \not = \underline{0} \;\;\text{ e }\;\; F(v) = \lambda v \,,\; \lambda \in \mathbb{R}
\] \(\lambda\) si dice \textbf{autovalore} dell'autovettore \(v\).

\textbf{Definizione:} Sia \(A \in M_n (\mathbb{R})\), diciamo che
\(\lambda\) è autovalore di \(A\) di \textbf{molteplicità algebrica}
\(m\) e scriviamo \[
\text{m}_a(\lambda) = m
\] se \(P_A(x) = (x- \lambda)^m g(x)\) e \(g(\lambda) \not = 0\), cioè
se \(m\) è la massimo potenza di \((x - \lambda)\) che divide \(P_A(x)\)
.

\textbf{Definizione:} Se \(\lambda\) è autovalore di \(A\) allora la
\textbf{molteplicità geometrica} di \(\lambda\) è \[
\text{m}_g(\lambda) = \dim(V_\lambda) = \dim(\text{Ker}(A - \lambda I))
\] Ovvero la dimensione dell'autospazio di \(\lambda\) .

\bigskip

\textbf{Proprietà equivalenti al fatto che un endomorfismo o una matrice
siano diagonalizzabili:}

Sia \(F : \mathbb{R}^n \to \mathbb{R}^n\) un'applicazione lineare. Le
seguenti proprietà sono equivalenti affinché \(F\) sia diagonalizzabile:

\begin{enumerate}
\def\labelenumi{\arabic{enumi}.}
\tightlist
\item
  Esiste una base \(\beta\) di \(\mathbb{R}^n\) tale che
  \(M^{{\beta}}_{\beta}(F)\) sia diagonale (definizione di endomorfismo
  diagonalizzabile)
\item
  Esiste una base di \(\mathbb{R}^n\) costituita da autovettori di \(F\)
\item
  Se \(\lambda_1, .., \lambda_k\) sono gli autovalori distinti di \(F\),
  con molteplicità geometriche \(n_1, .., n_k\), allora vale che
  \(n_1 + ... + n_k = n\)
\item
  Il polinomio caratteristico di
  \(A = M^{{\mathcal{C}}}_{\mathcal{C}}(F)\) si fattorizza in fattori
  lineari e, se \(\lambda_1, .., \lambda_k\) sono gli autovalori
  distinti di \(F\), si ha che
  \(\text{m}_a(\lambda_i) = \text{m}_g(\lambda_i)\) per ogni \(i\).
\end{enumerate}

\hypertarget{domanda-34}{%
\subsection{Domanda 34}\label{domanda-34}}

\begin{observationbox}[Domanda]
\begin{enumerate}
\def\labelenumi{\arabic{enumi}.}
\setcounter{enumi}{33}
\tightlist
\item
  Definizione di autovalore di un endomorfismo. Se un endomorfismo ha 0
  come autovalore, cosa possiamo dedurre e perché?
\end{enumerate}
\end{observationbox}

\textbf{Definizione:} Sia \(F: V \to V\) un endomorfismo.
\(\lambda \in \mathbb{R}\) è \textbf{autovalore} di \(F\) se esiste
\(v \in V, \;v \not = \underline{0}\) tale che \(F(v) = \lambda v\)

\bigskip

\textbf{Se un endomorfismo ha 0 come autovalore, cosa possiamo dedurre e
perché?}

Innanzitutto ricordiamo che il vettore nullo non può essere autovettore,
mentre 0 può essere autovalore. Nel caso in cui 0 sia autovalore, vuol
dire che per tutti i vettori \(v\) dell'autospazio \(V_0\) , vale che
\(F(v) = 0 \cdot v = \underline{0}\) e perciò \(V_0\) coincide con
\(\text{Ker}(F)\) (infatti inoltre
\(V_0 = \text{Ker}(A - 0I )= \text{Ker}(A)=\text{Ker}(F)\)) . Ma allora
vuol dire che \(\text{Ker}(F) \neq \{\underline{0}\}\), perciò \(F\) non
è iniettiva e perciò per il teoremone \(A\) non è invertibile,
\(\text{det}(A) = 0\) ecc\ldots{}

\hypertarget{domanda-35}{%
\subsection{Domanda 35}\label{domanda-35}}

\begin{observationbox}[Domanda]
\begin{enumerate}
\def\labelenumi{\arabic{enumi}.}
\setcounter{enumi}{34}
\tightlist
\item
  Definizione di prodotto scalare euclideo in \(\mathbb{R}^n\) e sue
  proprietà
\end{enumerate}
\end{observationbox}

\textbf{Definizione:} In \(\mathbb{R}^n\) il \textbf{prodotto scalare
euclideo} è una funzione: \[
\begin{aligned}
\langle \; , \; \rangle: \quad &\mathbb{R}^n \times \mathbb{R}^n &\to& \quad \mathbb{R} \\
&(\underline{x}, \underline{y}) &\to& \quad\langle \underline{x} , \underline{y} \rangle
\end{aligned}
\] tale che, se \(\underline{x} = (x_1, \dots, x_n)\) e
\(\underline{y} = (y_1, \dots, y_n)\): \[
\langle \underline{x} , \underline{y} \rangle = x_1y_1 + x_2y_2 + \dots + x_ny_n
\]

\textbf{Proprietà del prodotto scalare:}

\begin{enumerate}
\def\labelenumi{\arabic{enumi}.}
\tightlist
\item
  \(\langle x, y \rangle = \langle y, x \rangle\)
\item
  \(\langle x + y, z \rangle =\langle x, z \rangle +\langle y, z \rangle\)
\item
  \(\langle \lambda x, y \rangle = \lambda \langle x, y \rangle\)
\item
  \(\langle x, x \rangle \ge 0\)
\item
  \(\langle x, x \rangle = 0 \iff x = \underline{0}\)
\end{enumerate}

La proprietà (1) ci dice che il prodotto scalare è \textbf{simmetrico}.

Le proprietà (2) e (3) ci dicono che il prodotto scalare è
\textbf{bilineare}, cioè , se fisso \(z\), la funzione
\(x \to \langle x, z \rangle\) è una funzione
\(\mathbb{R}^n \to \mathbb{R}\) \textbf{lineare}.

Invece (4) e (5) dicono che il prodotto scalare euclideo è
\textbf{definito positivo}.

\hypertarget{domanda-36}{%
\subsection{Domanda 36}\label{domanda-36}}

\begin{observationbox}[Domanda]
\begin{enumerate}
\def\labelenumi{\arabic{enumi}.}
\setcounter{enumi}{35}
\tightlist
\item
  Definizione di sottospazio ortogonale di un sottospazio dato \(W\).
  Enunciare e dimostrare il seguente teorema: Se \(W\) è un sottospazio
  di \(\mathbb{R}^n\) di dimensione \(k\) allora \(W^\perp\) ha
  dimensione\ldots{} e \(W \cap W^{\perp} =\) \ldots{}
\end{enumerate}
\end{observationbox}

\textbf{Definizione:} Sia \(W \leq \mathbb{R}^n\). Definiamo il
\textbf{sottospazio ortogonale} \(W^{\perp}\): \[
W^{\perp} = \{ \underline{x} \in \mathbb{R}^n \mid \langle \underline{x}, w \rangle = 0, \forall \;w \in W\}
\]

\textbf{Teorema:} Se \(W\) è un sottospazio di \(\mathbb{R}^n\) di
dimensione \(k\) allora \(W^\perp\) ha dimensione \(n-k\) e
\(W \cap W^{\perp} = \{\underline{0}\}\) .

\textbf{Dimostrazione:}

Sia \(\beta = \{w_1, \dots, w_k\}\) base di \(W\). Allora \[
\underline{x} \in W^{\perp} \iff \langle \underline{x}, w_i \rangle = 0 \qquad\forall\, i = 1, \dots, k
\] Otteniamo così un sistema lineare di \(r\) equazioni in \(n\)
incognite con matrice associata: \[
\begin{pmatrix}  
(\quad w_1\quad ) &|\; 0 \\  
(\quad w_2\quad ) & |\;0 \\
\dots \\
(\quad w_k\quad ) &|\;0 \\
\end{pmatrix}  
= A \mid \underline{0}
\] Dopo averla ridotta a scala con Gauss abbiamo \(r(A) = k\), poiché
\(w_1, \dots, w_k\) sono linearmente indipendenti. Perciò le soluzioni
dipendono da \(n-k\) parametri. Quindi anche \(\dim(W)^{\perp} = n-k\)
poiché il sistema di equazioni ottenute non sono altro che le
\textbf{equazioni cartesiane di \(W^{\perp}\)} .

\bigskip

Dimostriamo ora che \(W \cap W^{\perp} = \{\underline{0}\}\).

Sia \(\underline{x} \in W \cap W^{\perp}\): \[
\underline{x} \in W^{\perp} \;\Longrightarrow\; \langle \underline{x}, w \rangle = 0 \quad\forall\, w \in W
\] e poi quindi \[
\underline{x} \in W \;\Longrightarrow\; \langle \underline{x}, \underline{x} \rangle = 0 
\] e per la positività del prodotto scalare \[
\underline{x} = \underline{0}
\]

\hypertarget{domanda-37}{%
\subsection{Domanda 37}\label{domanda-37}}

\begin{observationbox}[Domanda]
\begin{enumerate}
\def\labelenumi{\arabic{enumi}.}
\setcounter{enumi}{36}
\tightlist
\item
  Proiezione ortogonale di un vettore su di un sottospazio di
  \(\mathbb{R}^n\) .
\end{enumerate}
\end{observationbox}

\textbf{Definizione:} Sia \(W \leq \mathbb{R}^n\),
\(v \in \mathbb{R}^n\), la proiezione ortogonale di \(v\) su \(W\) è: \[
\text{proj}_W(v) = v_1
\] dove \(v_1\) è l'unico vettore tale che \(v= v_1 + v_2\) con
\(v_1 \in W\) e \(v_2 \in W^\perp\).

\bigskip

Infatti mostriamo ora che dato un vettore \(v \in W\), è sempre
possibile scrivere \(v\) come \(v = v_1 + v_2\) con \(v_1 \in W\) e
\(v_2 \in W^\perp\), e inoltre tale scrittura è unica.

\textbf{Dimostrazione:}

Sia \(\{w_1, ..., w_k\}\) base di \(W\) e sia \(\{u_1,\dots, u_{n-k}\}\)
base di \(W^\perp\). Consideriamo ora una base ortogonale
\(\{\widetilde{w}_1, ..., \widetilde{w}_k\}\) di \(W\) e una base
ortogonale \(\{\widetilde{u}_1, ..., \widetilde{u}_{n-k}\}\) di
\(W^{\perp}\).

Osserviamo che
\(\beta = \{\widetilde{w}_1,\dots, \widetilde{w}_k, \widetilde{u}_1, \dots, \widetilde{u}_{n-k}\}\)
è un insieme di vettori non nulli ortogonali tra loro, quindi sono
linearmente indipendenti. In tutto sono \(k + (n-k) = n\) quindi per GEL
\(\beta\) è base di \(\mathbb{R}^n\), quindi: \[
v = \lambda_1 \widetilde{w}_1 + ... + \lambda_k\widetilde{w}_k + \gamma_1 \widetilde{u}_1 + ... + \gamma_{n-k} \widetilde{u}_{n-k} = v_1 \in W + v_2 \in W^\perp
\]

Mostriamo ora l'unicità:

Se \(v = v_1 + v_2 = z_1 + z_2\), con \(v_1, z_1 \in W\) e
\(v_2, z_2 \in W^{\perp}\) , allora: \[
v_1 + v_2 = z_1 + z_2 \;\Longrightarrow\; v_1 - z_1 = z_2 - v_2
\] ma \(v_1 - z_1 \in W\) e \(z_2 - v_2 \in W^{\perp}\) e sappiamo che
\(W \cap W^{\perp} = \{\underline{0}\}\). Perciò: \[
v_1 = z_1 \qquad v_2 = z_2
\]

\bigskip

Vediamo ora come calcolare la proiezione ortogonale.

Sia \(\beta = \{w_1, ..., w_k, w_{k+1}, ..., w_n\}\) base
\textbf{ortonormale} di \(\mathbb{R}^n\) con \(w_1, ..., w_k \in W\) e
con \(w_{k+1}, ..., w_n \in W^{\perp}\).

Allora sia \[
v = \lambda_1 w_1 +...+ \lambda_k w_k + \lambda_{k+1}w_{k+1} + ... + \lambda_n w_n = v_1 \in W + v_2 \in W^\perp
\] Essendo \(\beta\) ortonormale allora
\(\lambda_i = \langle v, w_i \rangle\)

Quindi se abbiamo una base ortonormale di \(W\): \[
\text{proj}_W(v) = v_1 = \langle v, w_1 \rangle w_1 + \langle v, w_2 \rangle w_2 + \dots + \langle v, w_k \rangle w_k
\]

Se invece abbiamo solo una base \textbf{ortogonale}
\(\{z_1, ..., z_k\}\) la formula è: \[
\text{proj}_W(v) = \frac{\langle v, z_1 \rangle}{\langle z_1, z_1 \rangle}z_1 + \frac{\langle v, z_2 \rangle}{\langle z_2, z_2 \rangle}z_2 + \dots + \frac{\langle v, z_k \rangle}{\langle z_k, z_k \rangle}z_k
\]

\hypertarget{domanda-38}{%
\subsection{Domanda 38}\label{domanda-38}}

\begin{observationbox}[Domanda]
\begin{enumerate}
\def\labelenumi{\arabic{enumi}.}
\setcounter{enumi}{37}
\tightlist
\item
  Proprietà equivalenti affinché una matrice quadrata sia ortogonale.
\end{enumerate}
\end{observationbox}

Sia \(A \in M_n (\mathbb{R}^n)\). Sono equivalenti le seguenti
proprietà:

\begin{enumerate}
\def\labelenumi{\arabic{enumi}.}
\tightlist
\item
  \(\langle Au, Av \rangle = \langle u, v \rangle \quad \forall u, v \in \mathbb{R}^n\)
\item
  \(\langle Ax, Ax \rangle = \langle x, x \rangle \quad \forall x \in \mathbb{R}^n\)
\item
  \(A^T = A^{-1}\)
\item
  Le colonne (rispettivamente le righe) di \(A\) formano una base
  \textbf{ortonormale} di \(\mathbb{R}^n\)
\end{enumerate}

\textbf{Dimostrazione:}

\begin{itemize}
\item
  \begin{enumerate}
  \def\labelenumi{(\arabic{enumi})}
  \tightlist
  \item
    \(\Rightarrow\) (4) : Sia \(L_A : \mathbb{R}^n \to \mathbb{R}^n\)
    l'applicazione lineare tale che
    \(L_A (\underline{x}) = A\underline{x}\) . Allora : \[
    A = \left(  
    \begin{array}{cccc}  
    \big| & \big| &  \big| \\  
    L_A(e_1) 
    &  
    \cdots  
    &  
    L_A(e_1) \\  
    \big| & \big| & \big|  
    \end{array}  
    \right)
    = \left(  
    \begin{array}{cccc}  
    \big| & \big| &  \big| \\  
    Ae_1
    &  
    \cdots  
    &  
    Ae_n \\  
    \big| & \big| & \big|  
    \end{array}  
    \right)
    \] Per ipotesi
    \(\langle Ae_i, Ae_j \rangle = \langle e_i, e_j \rangle = \delta_{ij}\)
    . Cioè il prodotto scalare tra la colonna \(i\) e la colonna \(j\)
    di \(A\) è \(0\) se \(i \neq j\) e 1 se \(i = j\). Quindi le colonne
    di \(A\) sono ortonormali. Inoltre essendo vettori non nulli
    ortogonali sono indipendenti, e quindi base di \(\mathbb{R}^n\).
  \end{enumerate}
\item
  \begin{enumerate}
  \def\labelenumi{(\arabic{enumi})}
  \setcounter{enumi}{2}
  \tightlist
  \item
    \(\Rightarrow\) (4): Siano \(R_1, \dots, R_n\) le righe di \(A\) e
    quindi rispettivamente le colonne di \(A^T\). Sappiamo per ipotesi
    che \(A^T = A^{-1}\), quindi \(AA^T = I\) . Facendo il prodotto
    righe per colonne di \(AA^T\) ottengo che al posto \(ij\) c'è
    \(\langle R_i, R_j \rangle\) che d'altra parte è \(\delta_{ij}\) in
    quanto \(AA^T = I\). Perciò le righe di \(A\) sono una base
    ortonormale.
  \end{enumerate}
\item
  \begin{enumerate}
  \def\labelenumi{(\arabic{enumi})}
  \setcounter{enumi}{3}
  \tightlist
  \item
    \(\Rightarrow\) (3): Se sappiamo che
    \(\langle R_i, R_j \rangle = \delta_{ij}\) consideriamo il prodotto
    \(AA^T\) e scopriamo che viene \(I\) . Allora \(A^T = A^{-1}\)
  \end{enumerate}
\item
  \begin{enumerate}
  \def\labelenumi{(\arabic{enumi})}
  \setcounter{enumi}{3}
  \tightlist
  \item
    \(\iff\) (3): Con le colonne è analogo.
  \end{enumerate}
\end{itemize}

\hypertarget{domanda-39}{%
\subsection{Domanda 39}\label{domanda-39}}

\begin{observationbox}[Domanda]
\begin{enumerate}
\def\labelenumi{\arabic{enumi}.}
\setcounter{enumi}{38}
\tightlist
\item
  Enunciare il teorema spettrale.
\end{enumerate}
\end{observationbox}

\textbf{Teorema:} Sia \(F: \mathbb{R}^n \to \mathbb{R}^n\) un
endomorfismo simmetrico. Allora esiste una base ortonormale di
\(\mathbb{R}^n\) costituita da autovettori di \(F\) .

Ne seguono i seguenti corollari:

\begin{itemize}
\tightlist
\item
  Un endomorfismo simmetrico è sempre diagonalizzabile.
\item
  Se \(\lambda_1, \lambda_2\) sono autovalori distinti di \(F\) allora
  \(V_{\lambda_1}\) e \(V_{\lambda_2}\) sono ortogonali tra loro.
\item
  Se \(F = L_A\) simmetrico, allora esiste una matrice
  \textbf{ortogonale} \(P\) tale che \(P^TAP =D\) con \(D\) diagonale.
  Cioè \(A\) è \textbf{ortogonalmente simile} ad una matrice diagonale.
\end{itemize}

\hypertarget{domanda-40}{%
\subsection{Domanda 40}\label{domanda-40}}

\begin{observationbox}[Domanda]
\begin{enumerate}
\def\labelenumi{\arabic{enumi}.}
\setcounter{enumi}{39}
\tightlist
\item
  Definire le classi resto modulo \(n\), a partire da un'opportuna
  relazione di equivalenza in \(\mathbb{Z}\). In particolare
  caratterizzare \([a]_n = \{\dots\}\).
\end{enumerate}
\end{observationbox}

\textbf{Definizione:} Siano
\(a, b \in \mathbb{Z}, n \in \mathbb{N}, n \neq 0\) . Diciamo che \(a\)
è \textbf{congruo a \(b\) modulo \(n\)} e scriviamo \[
a \equiv_n b
\] se vale che \(n \mid (a-b)\) .

\bigskip

Inoltre la relazione di congruenza modulo \(n\) è una relazione di
equivalenza. Possiamo perciò definire le classi di congruenza o classi
di resto modulo \(n\):

\textbf{Definizione:} Definiamo \[
\left[a\right]_n = \{b \in \mathbb{Z} : b \equiv_n a\} = \{a + k \cdot n,\; k \in \mathbb{Z}\}
\] la \textbf{classe di congruenza} di \(a\), o \textbf{classe di resto
di \(a\) modulo \(n\)} .

\hypertarget{domanda-41}{%
\subsection{Domanda 41}\label{domanda-41}}

\begin{observationbox}[Domanda]
\begin{enumerate}
\def\labelenumi{\arabic{enumi}.}
\setcounter{enumi}{40}
\tightlist
\item
  Definire \(\mathbb{Z}_n\) e caratterizzarlo elencandone gli elementi,
  con dimostrazione.
\end{enumerate}
\end{observationbox}

\textbf{Definizione:} L'insieme quoziente di tutte le classi resto
modulo \(n\) è: \[
\mathbb{Z}_n = \{\left[a\right]_n : a \in \mathbb{Z}\}
\]

\textbf{Proposizione:} Vale che: \[
\mathbb{Z}_n = \{[0]_n, [1]_n, \dots, [n-1]_n \}
\] In particolare \(\mathbb{Z}_n\) contiene \(n\) elementi.

\textbf{Dimostrazione:}

Sia \(\left[a\right]_n\) una classe di equivalenza e mostriamo che
\(\left[a\right]_n\) è una delle \(n\) classi dell'elenco a destra.

Abbiamo \(a = n \cdot q + r, \; 0 \le r < n\) e abbiamo che
\(a \equiv_n r\) perché \(n \cdot q = a -r\) e quindi \(n \mid a-r\) .
Perciò: \[
a \equiv_n r \; \Longrightarrow \; \left[a\right]_n = [r]_n
\] e inoltre vale che: \[
r \in \{0, 1, \dots, n-1\}
\]

Mostriamo ora che le classi sono distinte, ovvero che se
\(i, j \in \{0, 1, \dots, n-1\}, \; i \neq j\) allora
\([i]_n \neq [j]_n\) .

Supponiamo che sia \([i]_n = [j]_n\) e sia \(i \ge j\), allora
\(i \equiv_n j\) quindi \(n \mid i - j\) , ma \(0 \leq i - j < n\) .
Quindi l'unica possibilità è che \(i - j = 0\), cioè \(i = j\) .

\hypertarget{domanda-42}{%
\subsection{Domanda 42}\label{domanda-42}}

\begin{observationbox}[Domanda]
\begin{enumerate}
\def\labelenumi{\arabic{enumi}.}
\setcounter{enumi}{41}
\tightlist
\item
  Enunciare e dimostrare il teorema: la classe \([a]_n\) è invertibile
  in \(\mathbb{Z}_n\) se e solo se\ldots{}
\end{enumerate}
\end{observationbox}

\textbf{Teorema:} la classe \([a]_n\) è invertibile in \(\mathbb{Z}_n\)
se e solo se \(\text{mcd}(n, a) = 1\)

\textbf{Dimostrazione:}

\(\Leftarrow\)) Sia \(\text{mcd}(a,n)=1\) , troviamo l'inverso di
\([a]_n\) .

Per il teorema di Bézout esistono \(u,v\in\mathbb Z\) tali che \[
1 = ua + vn
\] Passando alle classi modulo \(n\): \[
[1]_n = [ua+vn]_n
\] Per l'additività e la compatibilità del prodotto: \[
[1]_n
=
[ua]_n + [vn]_n
=
[u]_n[a]_n + [0]_n
=
[u]_n[a]_n
\] Quindi \[
[u]_n[a]_n=[1]_n
\] Pertanto \([u]_n\) è l'inverso di \([a]_n\) e quindi \([a]_n\) è
invertibile.

\bigskip

\(\Rightarrow\)) Sia \([a]_n\) invertibile e sia \([b]_n\) il suo
inverso. Allora \[
[a]_n[b]_n=[1]_n
\] Quindi \[
[ab]_n=[1]_n
\] ossia \[
ab \equiv_n 1 \quad \Longrightarrow \quad n \mid (ab-1)
\] Sia ora \[
d=\text{mcd}(a,n)
\] Poiché \(d\mid n\) e \(n\mid(ab-1)\), segue che \[
d\mid(ab-1)
\] Inoltre, essendo \(d\mid a\), si ha \[
d\mid ab
\]

Dunque \(d\) divide sia \(ab\) sia \(ab-1\).

Poiché un divisore comune divide anche la differenza, otteniamo \[
d \mid \bigl(ab-(ab-1)\bigr)
\] Ma \(ab-(ab-1)=1\), quindi \[
d\mid 1 \quad \Longrightarrow \quad d = 1
\] Pertanto \[
\text{mcd}(a,n)=1.
\]

\end{document}
